\section{Capítulo Teórico}\label{capuxedtulo-teuxf3rico}

\hyperref[capuxedtulo-teuxf3rico]{\textbf{Capítulo Teórico 1}}

\begin{quote}
\hyperref[teoria-socioluxf3gica-sobre-desindustrializauxe7uxe3o]{Teoria
Sociológica sobre desindustrialização 1}

\hyperref[desindustrializauxe7uxe3o-no-mundo]{Desindustrialização no
Mundo 1}

\hyperref[desindustrializauxe7uxe3o-no-brasil]{Desindustrialização no
Brasil 19}
\end{quote}

\textbf{Teoria Geral (15pg)}

Sociologia econômica contemporânea

\begin{quote}
principais linhas da soc econ

(redes, campos, instituições)

(organizações, mercados, grupos sociais e história)

A noção pervasiva de cultura

Predomínio de abordagens culturais no Brasil (Bourdieu)

A fuga para uma análise econômica não materialista e o espaço deixado
aos economistas

A opção por uma sociologia econômico-política como contraponto

Marx, o filho esquecido?

Crítica da crítica à sociologia ``estadocentrica'', ao marxismo em geral
e a sociologia do desenvolvimento (contexto dos anos 1980 em diante
-neoliberalismo, queda da URSS são razões para queda))

o poder e a violência não-simbólicos

China, Trump, desigualdade econômica e as justificativas para retomar
esse debate

Defesa de retomada estruturalista (marx) e centrada no estado

elementos centrais da teoria marxista ausentes no debate da soc
econômica

O problema do desenvolvimento econômico e das classes sociais (elites)

Estado e política como mediação do desenvolvimento econômico (questão do
poder político-econômico)

Necessidade de um contraponto institucionalista (polanyi → evans)

Institucionalismo

Sociologia econômica (histórica-comparativa e institucionalista) do
desenvolvimento
\end{quote}

institucionalismo (economico, cultural e político)

\begin{quote}
Diferentes abordagens mono institucionalista ou apenas relacionando
cultura e economia (marx, weber, schumpeter, veblen, elias)

Polanyi como institucionalismo econômico e político
\end{quote}

O caso do sul global (institucionalismo e dependência)

\begin{quote}
Imperialismo

Desenvolvimento e subdesenvolvimento

Teoria da dependência e desenvolvimento dependente

Institucionalismo e dependência (as instituições chave)
\end{quote}

\subsection{Teoria Sociológica sobre
desindustrialização}\label{teoria-socioluxf3gica-sobre-desindustrializauxe7uxe3o}

\subsection{Desindustrialização no
Mundo}\label{desindustrializauxe7uxe3o-no-mundo}

Introdução

Esse capítulo é uma revisão de literatura sobre processos de
desindustrialização. Seu objetivo é caracterizá-los em aspectos
sociológicos e históricos, identificando tanto suas características
generalizáveis como suas especificidades temporais e geográficas. Visa
também atualizar o leitor sobre os principais desenvolvimentos da
literatura de língua inglesa recente sobre o tema nas áreas de
sociologia e história.

O termo ``desindustrialização'' é comumente tratado como uma redução da
participação da manufatura e da indústria geral na economia
\href{https://www.zotero.org/google-docs/?SncwXs}{(Pula, 2017, p. 1)} ou
uma redução do emprego industrial em proporção ao estoque de empregos
\href{https://www.zotero.org/google-docs/?2eElZ8}{(Dandaneau, 2012, p.
1; Kollmeyer, 2020, p. 1; Lois González, 2019, p. 1)}. Em geral é
identificada pela mudança em indicadores de participação do setor
industrial no PIB
\href{https://www.zotero.org/google-docs/?UzkEfu}{(Rowthorn e Ramaswamy,
1999)} conjugadamente com indicadores de emprego no nível nacional
\href{https://www.zotero.org/google-docs/?L8IQDX}{(Tregenna, 2009)}.
Também pode ser abordado a partir de suas manifestações locais, como o
fechamento de fábricas
\href{https://www.zotero.org/google-docs/?8ZlqZh}{(Clarke, 2015; Hodson,
2019)}, regionais e setoriais
\href{https://www.zotero.org/google-docs/?hcZ5qK}{(Berger, Wicke e
Golombek, 2017; Durazzi e Geyer, 2020)}. Há na literatura também uma
crescente defesa de uma visão menos economicista do fenômeno, seja
abordando seus múltiplos significados para trabalhadores e comunidades
afetadas \href{https://www.zotero.org/google-docs/?XtB5a1}{(Cowie e
Heathcott, 2003)} ou tentando situá-lo como um processo histórico de
encerramento da sociedade industrial
\href{https://www.zotero.org/google-docs/?eZL9b3}{(Strangleman, 2017)}.

Neste trabalho propõe-se como definição de desindustrialização \textbf{a
transformação social associada à redução relativa da participação da
indústria na capacidade produtiva e no estoque de empregos em dado
território, que passa a se dar de forma acelerada e recorrente a partir
da segunda metade do século XX.} Essa mudança tem dimensões econômica,
política e cultural e é mediada por instituições tempo-espacialmente
enraizadas. Essa definição é resultado da revisão de literatura
apresentada, mas também da proposição de
\href{https://www.zotero.org/google-docs/?xdjSXH}{(Santos, 2016)} de
abordar o desenvolvimento econômico enquanto mudança social numa
perspectiva socioantropológica.

Enfatiza-se aqui, entretanto, seu caráter de transformação social
historicamente situada, como defendido em Strangleman
\href{https://www.zotero.org/google-docs/?whFleA}{(2017)} e Lawson
\href{https://www.zotero.org/google-docs/?LTSpDH}{(2020)}, e a
necessidade de conectar os estudos sobre trabalhadores e suas
comunidades às transformações de nível macro na economia política e na
divisão internacional do trabalho, como propõe High
\href{https://www.zotero.org/google-docs/?n0gdvx}{(High, 2023, p. 56)}.
Destaca-se a importância de compreender a aceleração dessa transformação
social, resultante da interação complexa de múltiplos fatores
contextuais, em vez de buscar causas únicas.

Entre as razões para investigar o fenômeno estão seus efeitos sobre e
relações com a política
\href{https://www.zotero.org/google-docs/?Ah5yEx}{(Iversen e Cusack,
2000)}, comércio
\href{https://www.zotero.org/google-docs/?MxvIWj}{(Alderson, 1999;
Alderson e Nielsen, 2002)}, emprego
\href{https://www.zotero.org/google-docs/?OizqO5}{(Beatty e Fothergill,
2020; Tomlinson, 2020)}, meio ambiente
\href{https://www.zotero.org/google-docs/?FAXcZI}{(Feltrin, Mah e Brown,
2022)} e seus impactos nas identidades dos indivíduos e comunidades
afetados por ele
\href{https://www.zotero.org/google-docs/?AiTsAR}{(Linkon, 2018; Mah,
2012; Nayak, 2006; Perchard, 2013)}. Inicialmente identificada com o
rust belt americano e regiões europeias como o Ruhr, Alemanha, ou norte
da Inglaterra \href{https://www.zotero.org/google-docs/?0FGLrz}{(Emery,
2019; High, 2013; Strangleman, Rhodes e Linkon, 2013)}, a discussão tem
ganhado crescente espaço fora do Norte global, sendo seu aparecimento
``prematuro'' visto como um desafio para países de industrialização
tardia \href{https://www.zotero.org/google-docs/?NLLOOu}{(Rodrik,
2016)}.

No Brasil, as discussões centram-se sobre sua ocorrência, causas e
consequências \href{https://www.zotero.org/google-docs/?YVvSKG}{(Maia,
2020; Marconi e Rocha, 2012; Morceiro e Guilhoto, 2020; Nassif, 2008;
Oreiro e Feijó, 2010)}, mas o debate acaba centrado no campo da
economia, limitando a compreensão da complexidade do fenômeno. Como
exceção, Santos et al
\href{https://www.zotero.org/google-docs/?lvTNTQ}{(Santos \emph{et al.},
2024)} apontam sua relação com mudanças na dinâmica eleitoral
brasileira. Por sua vez, políticas públicas como a ``Nova Indústria
Brasil'' partem da premissa da existência da desindustrialização e visam
revertê-la ou amenizá-la como parte de uma estratégia de desenvolvimento
econômico \href{https://www.zotero.org/google-docs/?PCangB}{(CNDI e
MDIC, 2024, p. 5)}

O texto a seguir está dividido em quatro partes que abordam: 1 -
critérios de seleção e características da revisão de literatura; 2 -
características econômicas, culturais e políticas da
desindustrialização; 3- a temporalidade, lugares e grupos sociais
privilegiados no seu debate; 4 - suas causas, consequências e relações
territoriais e setoriais; 5- discussão e conclusões.

1. Características da revisão de literatura

A revisão de literatura realizada é exploratória e crítica
\href{https://www.zotero.org/google-docs/?2AtMYd}{(Jesson, Matheson e
Lacey, 2011, p. 15)}, incorporando elementos de uma revisão sistemática
\href{https://www.zotero.org/google-docs/?WshOkn}{(Efron e Ravid, 2019,
p. 25)}, mas com flexibilização da cobertura do material. Foram
utilizados termos de busca na \emph{Web of
Science}\footnote{Os termos de busca também são denominados comumente
  como \emph{querys} e tipicamente permitem o encadeamento booleano de
  combinações de termos que devem estar presentes ou serem excluídos dos
  resultados.}\footnote{A \emph{query} utilizada teve o seguinte
  formato: \emph{(deindustrialization OR deindustrialisation OR
  de-industrialization OR de-industrialisation OR reindustrialization OR
  reindustrialisation OR re-industrialization OR re-industrialisation OR
  ``manufacturing decline'' OR ``decline in manufacturing'' OR
  ``industrial closings'' OR ``shift to service economy'' OR
  ``transition to service economy'')}, sendo os termos procurados no
  título, resumo ou palavras chave dos artigos (critério ``topics'' pela
  nomenclatura da \emph{Web of Science}).} para encontrar 205 artigos em
inglês nas áreas de sociologia e história, publicados entre 2014 e 2023
e revisados por pares.

Esses artigos foram ordenados por número de citações e avaliados quanto
ao seu tópico central, abordagem do conceito de desindustrialização,
dimensões consideradas e metodologia (cf. tabela 11 abaixo). Vinte e
três artigos foram selecionados com base nesses critérios. Em seguida,
foram extraídas as referências citadas pelo conjunto de 205 artigos,
resultando em uma lista das fontes bibliográficas mais utilizadas nos
últimos dez anos. Foram selecionados os artigos revisados por pares mais
relevantes, adicionando assim vinte artigos e os cinco livros mais
citados ao \emph{corpus}. No total, tem-se a seleção sistemática de 43
artigos e cinco livros como base para a revisão de literatura. Por fim,
critérios não sistemáticos foram utilizados para preencher lacunas
específicas relativas ao caso brasileiro e contextualizar a discussão na
introdução acima.

Tabela 11: critérios de priorização de artigos

\begin{longtable}[]{@{}
  >{\raggedright\arraybackslash}p{(\linewidth - 2\tabcolsep) * \real{0.1538}}
  >{\raggedright\arraybackslash}p{(\linewidth - 2\tabcolsep) * \real{0.8462}}@{}}
\toprule\noalign{}
\begin{minipage}[b]{\linewidth}\raggedright
Critério
\end{minipage} & \begin{minipage}[b]{\linewidth}\raggedright
Descrição da avaliação
\end{minipage} \\
\begin{minipage}[b]{\linewidth}\raggedright
Tópico Central
\end{minipage} & \begin{minipage}[b]{\linewidth}\raggedright
(+) A desindustrialização é um tópico central do artigo? (-) Há outras
temáticas que dividem o foco do artigo?
\end{minipage} \\
\begin{minipage}[b]{\linewidth}\raggedright
Abordagem do conceito
\end{minipage} & \begin{minipage}[b]{\linewidth}\raggedright
(+) Há uma definição clara do conceito? (+) O conceito é considerado em
sua complexidade de causas e consequências? (-) É tratado como apenas
variável causal para outro fenômeno?
\end{minipage} \\
\begin{minipage}[b]{\linewidth}\raggedright
Dimensões consideradas
\end{minipage} & \begin{minipage}[b]{\linewidth}\raggedright
O artigo aborda a desindustrialização em dimensões: (+) econômica, (+)
social/cultural e (+) política?
\end{minipage} \\
\begin{minipage}[b]{\linewidth}\raggedright
Metodologia
\end{minipage} & \begin{minipage}[b]{\linewidth}\raggedright
(+) Artigo possui descrição de metodologia clara? (+) Descreve a fonte
dos dados analisados?
\end{minipage} \\
\multicolumn{2}{@{}>{\raggedright\arraybackslash}p{(\linewidth - 2\tabcolsep) * \real{1.0000} + 2\tabcolsep}@{}}{%
\begin{minipage}[b]{\linewidth}\raggedright
Nota: (+) indica aspecto positivo; (-) indica aspecto negativo;
pontuação de 0, 0.5 ou 1 em cada critério; selecionados artigos que
pontuaram 3 ou mais
\end{minipage}} \\
\midrule\noalign{}
\endhead
\bottomrule\noalign{}
\endlastfoot
\end{longtable}

Fonte: autoria própria

Realizando uma inferência a partir da revista de publicação, ou da
formação acadêmica do autor, pode-se agrupar os artigos da seguinte
maneira: sociologia (15), história (18), economia (4), ciência política
(3), estudos culturais (1), \emph{social policy} (1) e geografia (1). Os
livros, a partir da formação dos autores, podem ser classificados em:
sociologia (1), história (2), economia (1) e estudos de literatura (1).
Tal agrupamento, entretanto, não dá conta completa da natureza
interdisciplinar dos trabalhos do campo. Tal interdisciplinaridade se
reflete na heterogeneidade das linhas teóricas utilizadas, estando
presentes abordagens neoclássicas
\href{https://www.zotero.org/google-docs/?kkP1z2}{(Scheiring e King,
2023)}, marxistas
\href{https://www.zotero.org/google-docs/?t9NhX6}{(High, 2013, p. 997)},
institucionalistas
\href{https://www.zotero.org/google-docs/?VlDdus}{(Lee, 2016)}, da
história econômica
\href{https://www.zotero.org/google-docs/?TRqrV5}{(Van Nederveen
Meerkerk, 2017)} e da sociologia da estratificação e mobilidade social
\href{https://www.zotero.org/google-docs/?MEPAJj}{(Kollmeyer, 2018)},
entre outras.

Do ponto de vista metodológico, predominam na literatura trabalhos
primordialmente qualitativos. Há forte presença de trabalhos que
utilizam metodologias de história oral
\href{https://www.zotero.org/google-docs/?cIWlm8}{(Clark e Gibbs, 2020;
Clarke, 2015; Hodson, 2019)}\hl{, entrevistas}
\href{https://www.zotero.org/google-docs/?u1rytv}{(Hoekstra, 2020;
Nettleingham, 2019; Scheiring, 2020; Scheiring e King, 2023)} \hl{e em
menor grau análises documentais}
\href{https://www.zotero.org/google-docs/?kgYUU2}{(Phillips, 2015;
Phillips, Wright e Tomlinson, 2019)}\hl{, etnografias}
\href{https://www.zotero.org/google-docs/?yB9B2v}{(Lewis, 2016)} \hl{e
análises institucionalistas comparativas}
\href{https://www.zotero.org/google-docs/?rZVpjC}{(Lee, 2016)}\hl{.
Alguns trabalhos empregam métodos mistos, combinando análises históricas
com análises de dados quantitativos}
\href{https://www.zotero.org/google-docs/?CYu4MG}{(Tomlinson, 2020; Van
Nederveen Meerkerk, 2017)}\hl{.}

\hl{Já trabalhos de ênfase primordialmente quantitativa são minoria, se
destacando preocupações com o impacto econômico da desindustrialização
na desigualdade econômica e mobilidade social}
\href{https://www.zotero.org/google-docs/?hVR5XO}{(Berger e Engzell,
2022; Kollmeyer, 2018)} \hl{ou suas causas do ponto de vista econômico}
\href{https://www.zotero.org/google-docs/?odDWC6}{(Kollmeyer, 2018;
Rodrik, 2016)}\hl{. Estes são na maioria baseados em técnicas de
regressão multivariada. Por fim, alguns trabalhos têm preocupação
primordialmente teórica}
\href{https://www.zotero.org/google-docs/?aGjtPq}{(Strangleman, 2017;
Tomlinson, 2016)} \hl{enquanto outros são revisões de literatura}
\href{https://www.zotero.org/google-docs/?lTQk3H}{(Emery, 2019; High,
2013; Lawson, 2020; Strangleman e Rhodes, 2014; Strangleman, Rhodes e
Linkon, 2013)}\hl{.}

2. Características econômicas, culturais e políticas do processo de
desindustrialização

Diversas maneiras pelas quais o processo de desindustrialização é
compreendido na literatura podem ser agregadas a partir da ênfase em
seus aspectos econômicos, culturais ou políticos. Obviamente a maioria
dos trabalhos não adere estritamente a uma das categorias, em geral
sendo valorizada a interação de duas ou mais dessas categorias.

No que concerne à dimensão econômica, a redução do emprego industrial é
o principal indicador do fenômeno na literatura examinada
\href{https://www.zotero.org/google-docs/?dfS8v8}{(Beatty e Fothergill,
2020; Berger e Engzell, 2022; Clark e Gibbs, 2020, p. 42; Phillips,
Wright e Tomlinson, 2019; Scheiring, 2020; Tomlinson, 2016, p. 84; Van
Nederveen Meerkerk, 2017)}, em diversos casos analisada a partir da
experiência de fechamento de fábricas
\href{https://www.zotero.org/google-docs/?PH688f}{(Clarke, 2015; Hodson,
2019; Lewis, 2016)}, saída de empresas de regiões
\href{https://www.zotero.org/google-docs/?aXSL6s}{(Highsmith, 2014)} ou
declínio de setores industriais particulares, como carvão e aço
\href{https://www.zotero.org/google-docs/?AU5O44}{(Berger, Wicke e
Golombek, 2017)}. Alguns autores tratam o processo, de forma menos
específica, como uma redução da capacidade produtiva
\href{https://www.zotero.org/google-docs/?AHGVw8}{(Bluestone e Harrison,
1982)}, um encolhimento dos setores industriais centrais
\href{https://www.zotero.org/google-docs/?kdGwX7}{(Durazzi e Geyer,
2020)} ou uma redução da produção -- \emph{output} no original --
industrial \href{https://www.zotero.org/google-docs/?l64MF9}{(Van
Nederveen Meerkerk, 2017, p. 1219--1220)}.

Para Bluestone e Harrison
\href{https://www.zotero.org/google-docs/?CdGPle}{(1982)}, Kollmeyer
\href{https://www.zotero.org/google-docs/?goGCB5}{(2018, p. 1)} e Gibbs
e Clark \href{https://www.zotero.org/google-docs/?jifGFI}{(2020, p. 40)}
a contração do setor industrial está relacionada à crescente mobilidade
de capital e para Tomlinson
\href{https://www.zotero.org/google-docs/?mcJpA6}{(2016, 2020, p. 199)}
ao crescimento do setor terciário e transição para uma economia de
serviços. Tal transição é associada à crescente terceirização e aumento
do trabalho precário em Lee
\href{https://www.zotero.org/google-docs/?M0nGdv}{(2016, p. 79)} e a
salários mais baixos em Winant
\href{https://www.zotero.org/google-docs/?TV45XC}{(2019, p. 107)}. Ainda
sobre o mercado de trabalho, há ênfase numa crescente dualização, a
separação deste em camadas que mantêm certo grau de proteção social e
uma crescente camada de trabalhadores submetidos a piores condições de
trabalho e menores salários
\href{https://www.zotero.org/google-docs/?Z0WlY5}{(Durazzi e Geyer,
2020, p. 120; Lee, 2016, p. 86)}.

\hl{A conceituação do fenômeno, porém, vai além, incorporando dimensões
culturais e políticas. É notável a influência programática de Cowie e
Heathcott nos estudos de desindustrialização ao propor repensar sua
``cronologia, memória, relações espaciais, cultura e políticas'' {[}e
analisar{]} ``os múltiplos significados de uma das mais importantes
transformações do século XX''}
\href{https://www.zotero.org/google-docs/?OCJTFh}{(Cowie e Heathcott,
2003, p. 2)}\footnote{Tradução própria}\hl{.}

\hl{Dentro da dimensão cultural, algumas importantes considerações são
sobre como essa mudança representa uma transição de um modo de vida e
cujos impactos se estendem longamente. Linkon
\href{https://www.zotero.org/google-docs/?NzkGTV}{(2018)} vai enfatizar
essa continuidade pela expressão "meia vida da desindustrialização",
destacando os efeitos de longo prazo sobre trabalhadores e comunidades
afetados pelos fechamentos de fábricas. Mah
\href{https://www.zotero.org/google-docs/?tI9fox}{(2012)} irá olhar o
constante declínio dessas comunidades sob a perspectiva de uma
"ruinação", enfatizando assim seu aspecto de processo.}

\hl{Já Strangleman
\href{https://www.zotero.org/google-docs/?AYbZdZ}{(2017)} avança a
discussão propondo a leitura dessa transformação como um momento
histórico simétrico à industrialização. Em sua leitura de Thompson}
\href{https://www.zotero.org/google-docs/?iE7wh2}{\hl{(1961.}
\emph{\hl{apud}} \hl{Strangleman, 2017)}} \hl{ele enfatizará este
momento como uma transição que carrega consigo as marcas do período
anterior, um fechamento da era industrial
\href{https://www.zotero.org/google-docs/?IjVSYR}{(Strangleman, 2017, p.
467)}. De forma similar, Hoekstra
\href{https://www.zotero.org/google-docs/?KIdYBO}{(2020, p. 694)}
encapsula as discussões sobre transformações relativas às comunidades
como constituindo a perda de uma forma de vida. Scheiring e King
\href{https://www.zotero.org/google-docs/?9kbgx1}{(Scheiring e King,
2023, p. 144)}vão propor ver a indústria como uma instituição social
cuja sua dissolução provoca uma desintegração com amplas implicações
para o tecido social.}

\hl{Memória e nostalgia são também mobilizados}
\href{https://www.zotero.org/google-docs/?Dc2GGA}{(Clark e Gibbs, 2020;
Clarke, 2015; Lewis, 2016; Nettleingham, 2019)}\hl{. Porém, longe de
estarem ligados meramente ao passado, esses conceitos ressaltam como os
atores sociais mobilizam as narrativas do passado no presente enquanto
recursos, possibilitando assim a transformação do futuro. Dizem respeito
menos ao passado e mais aos ``presentes alternativos''}
\href{https://www.zotero.org/google-docs/?clyMVg}{(Clarke, 2015, p.
119--123)} \hl{possíveis de existir e aos futuros possíveis de serem
imaginados.}

Já sobre a dimensão política, de conflito e coordenação, \hl{Phillips}
\href{https://www.zotero.org/google-docs/?7iYBD0}{(2015, p. 551)}
\hl{destaca o caráter desejado e politizado da desindustrialização.
Mostra, pela história do fechamento de uma mina de carvão, como as
decisões -- e propensões -- dos atores políticos e das organizações
econômicas são baseadas em suas relações sociais e tensões e que o
processo sempre pode ser gerido de formas distintas}
\href{https://www.zotero.org/google-docs/?T6n73h}{(Phillips, 2015, p.
571--572)}\hl{. Já Clarke}
\href{https://www.zotero.org/google-docs/?YgWzzI}{(Clarke, 2015, p.
108)} \hl{mostra como a desindustrialização está imbricada numa
narrativa de ``modernização'' pela qual o fenômeno é considerado
inevitável, funcionando como estratégia por parte de representações
oficias do poder público de forma a despolitizar as discussões sobre
projetos relativos aos antigos territórios industriais.}

\hl{Por sua vez, Clark e Gibbs}
\href{https://www.zotero.org/google-docs/?QpVkIs}{(2020, p. 45)}
\hl{apontam os desequilíbrios de poder e distintas experiências de
atores da sociedade civil, econômicos e políticos em regiões passando
por entradas ou saídas de investimentos. Já numa dimensão mais ampla,
Tomlinson} \href{https://www.zotero.org/google-docs/?FkA0zT}{(2016)}
\hl{aborda a desindustrialização como uma metanarrativa explicativa das
transformações da Grã-Bretanha no século XX e destaca suas consequências
para a economia política e bem-estar econômico}
\href{https://www.zotero.org/google-docs/?80xC3Z}{(Tomlinson, 2016, p.
86)}\hl{. Como um todo, há assim uma noção de desindustrialização como
uma disputa pertencente à economia política na qual claramente alguns
atores -- do mercado e Estado -- possuem posição privilegiada em relação
aos trabalhadores e suas comunidades.}

\hl{Por fim}, a gestão da desindustrialização no tempo presente é
tratada mais centralmente em dois conceitos político-culturais. Rhodes
\href{https://www.zotero.org/google-docs/?dYyfpM}{(2013, p. 64)} e Byrne
\href{https://www.zotero.org/google-docs/?izvyTn}{(2002)} mobilizam a
ideia de estrutura de sentimentos, de Raymond Williams, como forma de
descrição da maneira como o processo histórico é percebido no dia a dia
por parte daqueles envolvidos. Já a expressão ``economia moral'', em
Gibbs \href{https://www.zotero.org/google-docs/?7BP6gP}{(2018)} e
Tomlinson \href{https://www.zotero.org/google-docs/?ULOwso}{(2011)}, em
grande parte referenciando Thompson
\href{https://www.zotero.org/google-docs/?cOHBIG}{(1966)}, diz respeito
a aspectos normativos, difusos nas comunidades, sobre como a gestão da
transformação econômica deveria se dar, especialmente do ponto de vista
da ação estatal e da mitigação, ou não, dos impactos sobre trabalhadores
e comunidades.

3. Lugares, periodização e personagens

\hl{A maior parte das análises concentra-se em países com alto grau de
renda \emph{per capita} e industrialização antiga, em especial na Europa
e América do Norte. As mais frequentes remetem à Grã-Bretanha, aparente
reflexo de sua importância histórica no processo de industrialização.
Tomlinson} \href{https://www.zotero.org/google-docs/?PSiiOG}{(2016)}
\hl{toma como recorte a Grã-Bretanha como um todo a partir de 1945,
Nettleingham} \href{https://www.zotero.org/google-docs/?DzKgjo}{(2019)}
\hl{dá ênfase à cidade inglesa de Faversham, enfatizando os precedentes
históricos dos fechamentos em massa dos anos 80 e Beatty e Fothergill}
\href{https://www.zotero.org/google-docs/?iHJANV}{(2020)} \hl{expandem o
escopo para um conjunto de antigas cidades industriais inglesas. Já
Lewis} \href{https://www.zotero.org/google-docs/?cZ7T4H}{(2016)}
\hl{trata em específico da região leste de Manchester. Aparecem ainda
discussões sobre a Irlanda do Norte}
\href{https://www.zotero.org/google-docs/?adDhwJ}{(Hodson, 2019)} \hl{e
Escócia} \href{https://www.zotero.org/google-docs/?7OesxW}{(Perchard,
2013; Phillips, 2013, 2015; Phillips, Wright e Tomlinson, 2019)}\hl{.
Há, ainda, debates centrados na Europa continental, como na região do
Ruhr, na Alemanha}
\href{https://www.zotero.org/google-docs/?XOxNcJ}{(Berger, Wicke e
Golombek, 2017)}\hl{, França}
\href{https://www.zotero.org/google-docs/?SSjZ1w}{(Clarke, 2015)} \hl{e
Hungria pós transição para uma economia de mercado}
\href{https://www.zotero.org/google-docs/?i9n36z}{(Scheiring, 2020;
Scheiring e King, 2023)}\hl{.}

\hl{Também central é o caso dos EUA, tratado na escala nacional}
\href{https://www.zotero.org/google-docs/?JJF36r}{(Berger e Engzell,
2022; Kollmeyer, 2018)}\hl{, regional, especialmente quanto ao que se
convencionou chamar de ``\emph{rust belt''} americano}
\href{https://www.zotero.org/google-docs/?ZmBFJG}{(Bluestone e Harrison,
1982; High, 2003)} \hl{e local, com particular atenção a cidades como
Flint} \href{https://www.zotero.org/google-docs/?t9HVJz}{(Highsmith,
2014)} \hl{e Pittsburgh}
\href{https://www.zotero.org/google-docs/?SRxFGq}{(Winant, 2019)}\hl{.
Ainda na América do Norte, Clark e Gibbs}
\href{https://www.zotero.org/google-docs/?9HXi1x}{(2020)} \hl{enfocam a
Nova Escócia, Canadá, ecoando a presença do país como um dos pioneiros
nos estudos de desindustrialização}
\href{https://www.zotero.org/google-docs/?h6soGG}{(High, 2013, p.
994)}\hl{.}

\hl{Alguns trabalhos destacam-se pelo caráter comparativo internacional,
como: Mah} \href{https://www.zotero.org/google-docs/?qWOBi4}{(2012)}
\hl{que traz comparações de cidades americana, inglesa e russa; Durazzi
e Geyer} \href{https://www.zotero.org/google-docs/?7BvhKX}{(2020)}\hl{,
centrando a discussão nos contextos Austríaco e Alemão; e Lee}
\href{https://www.zotero.org/google-docs/?2GMa3o}{(2016)}\hl{, que
analisa comparativamente Japão, Coreia do Sul e Taiwan. Este último se
mostra uma das poucas exceções a fugir do eixo Norte Atlântico.}

\hl{Do ponto de vista da escala analítica, a maior parte dos artigos
trabalha com a escala do Estado-Nação ou país. Também são frequentes
artigos tomando como escala analítica determinada cidade em específico,
às vezes expandindo sua análise para a região metropolitana, como em
Highsmith}
\href{https://www.zotero.org/google-docs/?GOOAao}{(2014)}\hl{. Em alguns
exemplos mais raros há ênfase regional, como em Berger; Wicke e
Golombek} \href{https://www.zotero.org/google-docs/?WBJ1XC}{(2017)}
\hl{sobre a região do Ruhr ou Scheiring}
\href{https://www.zotero.org/google-docs/?hUQSGX}{(2020)}\hl{, o qual
seleciona um conjunto de cidades húngaras tipicamente industriais antes
da transição para economia de mercado. Alguns trabalhos deslocam a
escala ainda mais para a unidade industrial, como Hodson}
\href{https://www.zotero.org/google-docs/?bh9WW6}{(2019)}\hl{, no qual
há ênfase no encerramento de um icônico estaleiro e Phillips}
\href{https://www.zotero.org/google-docs/?nAwPpr}{(2015)}\hl{, onde se
descreve em detalhe o fechamento de uma mina de carvão, ambos na
Escócia. Dada a natureza multiescalar dos processos econômicos, uma
lacuna dessa literatura reside na falta de trabalhos operando em
diversas escalas analíticas simultaneamente ou com abordagens de rede.}

\hl{É Importante salientar como a escolha da escala analítica muitas
vezes é determinante na identificação da existência de processos de
desindustrialização, em especial em situações que envolvem realocação
produtiva. Um dos melhores exemplos é como Bluestone e Harrison
\href{https://www.zotero.org/google-docs/?fH32OW}{(1982)} enfatizam
simultaneamente a ocorrência de desindustrialização no norte dos EUA
enquanto observam a industrialização no sul do país. Também pioneiro, e
poucas vezes repetido, é o modo como esses mesmos autores abordam os
impactos sociais dos dois processos, mostrando-os como duas faces de um
mesmo fenômeno relacionado à crescente mobilidade de capital.}

\hl{Essa transformação, entretanto, não segue um ritmo linear ou
simultâneo em todos os territórios. Considerando a escala nacional para
fins didáticos, as análises do caso americano apontam um intenso período
de desinvestimento produtivo nos anos 60 até os anos 90}
\href{https://www.zotero.org/google-docs/?Wbf6Ja}{(Bluestone e Harrison,
1982)}\hl{, com intensificação nos anos 70}
\href{https://www.zotero.org/google-docs/?epVUGe}{(Highsmith, 2014, p.
42; Strangleman, 2017, p. 474; Winant, 2019)}\hl{. Kollmeyer}
\href{https://www.zotero.org/google-docs/?00gScK}{(2009)} \hl{chega a
pontuar seus sinais já no pós Segunda Guerra Mundial, enquanto Highsmith
\href{https://www.zotero.org/google-docs/?BapSZJ}{(2014)} situa-a como
centrada nos anos 1960 em diante, porém ligado intimamente a
descentralização do investimento produtivo já a partir dos anos 1940.
Sobre o Canadá, Clark e Gibbs}
\href{https://www.zotero.org/google-docs/?840iJG}{(2020)} \hl{destacam
um processo a partir dos anos 1970 em diante, intensificado nos anos
1980.}

\hl{Os casos da Grã-Bretanha e europeus continentais exibem padrão
similar. No primeiro, há evidências iniciais de desindustrialização
cujas causas remontam ao período prévio a 1945}
\href{https://www.zotero.org/google-docs/?a0K08l}{(Tomlinson,
2020)}\hl{, mas cujos desenvolvimentos passam a ser claramente vistos em
fins dos anos 1950, abrangendo a Inglaterra, a Escócia e a Irlanda do
Norte e se aceleram ao longo dos anos 1960 e 1970}
\href{https://www.zotero.org/google-docs/?NhrNAK}{(Hodson, 2019;
Phillips, 2015; Tomlinson, 2016)}\hl{. O período de 1979 a 1982 é,
particularmente, crítico para a Grã-Bretanha como aponta Tomlinson}
\href{https://www.zotero.org/google-docs/?J7dR83}{(2016, p. 87)}\hl{, de
1975 a 1981 para a Escócia}
\href{https://www.zotero.org/google-docs/?QXPNE6}{(Phillips, Wright e
Tomlinson, 2019)}\hl{, se mantendo intenso ao longo dos anos 1980}
\href{https://www.zotero.org/google-docs/?DcU6gD}{(Beatty e Fothergill,
2020, p. 2)}\hl{.}

\hl{Tanto EUA como Inglaterra, demonstram, então, alguns indícios de uma
intensificação da desindustrialização e seus impactos com a eleição dos
governos Reagan (1980) e Tatcher (1979) respectivamente, mas há uma
avaliação de que está é prévia a esses governos e já vinha se acelerando
mesmo antes da ascensão de governos classificados como neoliberais}
\href{https://www.zotero.org/google-docs/?fFlzMu}{(Tomlinson,
2016)}\hl{. Winant argumenta inclusive sobre como os incrementos de
custos relativos à seguridade social decorrentes da desindustrialização
foram na verdade elemento importante nas disputas da economia política
que levaram à ascensão de tais governos}
\href{https://www.zotero.org/google-docs/?mfJsD7}{(Winant, 2019, p.
113--114)}\hl{.}

\hl{Indo em direção a Europa continental, na Alemanha duas ondas se
sobrepõem: uma ligada à indústria de extração de carvão nos anos 1950 em
diante e outra à siderurgia, dos anos 1970 em diante na região do Ruhr}
\href{https://www.zotero.org/google-docs/?h9qam5}{(Berger, Wicke e
Golombek, 2017)}\hl{. No entanto, o território exibiu maior gradualidade
no fechamento de fábricas, prolongando o processo e mitigando seus
custos sociais}
\href{https://www.zotero.org/google-docs/?MVk663}{(Berger, Wicke e
Golombek, 2017)}\hl{. Já na França, é possível ver um sinal de como a
aceleração maior deu-se apenas em fins dos anos 1980 e início dos anos
1990} \href{https://www.zotero.org/google-docs/?CGshOg}{(Clarke,
2015)}\hl{.}

\hl{O Leste Europeu e o Leste Asiático, por sua vez, apresentam uma
periodização distinta. Na Hungria esse processo se deu de forma
acelerada ao longo dos anos 1990, possivelmente contribuindo para uma
intensificação de seus impactos}
\href{https://www.zotero.org/google-docs/?O6VLTx}{(Scheiring, 2020;
Scheiring e King, 2023)}\hl{. Já no Leste Asiático, há identificação de
desindustrialização dos anos 1990 em diante, porém, sua ocorrência é
mais gradual e mitigada pela ação estatal}
\href{https://www.zotero.org/google-docs/?8gOF1N}{(Lee, 2016)}\hl{. Há
notável ausência, nessa literatura, de discussões sobre o caso Chinês.}

\hl{Em resumo, apesar de não significar a inexistência de casos
temporalmente anteriores, como mostram outros trabalhos
\href{https://www.zotero.org/google-docs/?fMnUVP}{(Clingingsmith e
Williamson, 2008; Pamuk e Williamson, 2011)}, há um consenso na
literatura sobre sua maior intensidade ocorrendo entre os anos 1960 e
fins dos anos 1980, especialmente na Europa Ocidental e EUA. Territórios
caracterizados por uma industrialização mais tardia}\footnote{Desindustrialização
  mais tardia aqui refere-se aos países onde seu processo de
  industrialização ocorreu durante o século XX, em contraposição a
  processos de desindustrialização nos séculos XVIII e XIX
  \href{https://www.zotero.org/google-docs/?MTwFmM}{(Mello, 1982)}}
\hl{ou com nível de renda \emph{per capita} menor que os casos americano
e europeu parecem ter seus processos situados a partir dos anos 1980,
com intensificação ao longo dos anos 1990 e 2000, dialogando com o
diagnóstico de Rodrik
\href{https://www.zotero.org/google-docs/?YSyJ0m}{(2016)}.}

Da mesma forma, algumas categorias de indústria são privilegiadas no
debate, particularmente em função de sua relevância para os territórios
enfatizados. O exemplo alemão, como já apresentado, mostra um processo
de desindustrialização iniciado na indústria extrativa, de
carvão,\footnote{Vale pontuar a relação entre a formação da União
  Europeia e a gestão dos recursos de carvão e aço no continente, sendo
  a primeira resultado direto da formação da Comunidade Europeia do
  Carvão e do Aço (CECA) em 1951
  \href{https://www.zotero.org/google-docs/?L9vJ4c}{(Valls, 2016)}} que
posteriormente se estende para a indústria de transformação pesada, a
siderurgia \href{https://www.zotero.org/google-docs/?cCMYRz}{(Berger,
Wicke e Golombek, 2017)}. A indústria de carvão também é a indústria
central examinada no exemplo Britânico
\href{https://www.zotero.org/google-docs/?7lhlhr}{(Tomlinson, 2016)}. Há
uma lacuna nos estudos a respeito da composição típica da indústria em
cada território e seus impactos. Assim, não pode ser excluída a
possibilidade de fatores de caráter subsetorial influenciarem as
trajetórias observadas.

Já nos EUA e no leste asiático predominam estudos sobre o fechamento de
fábricas e unidades produtivas ligadas à manufatura. No \emph{rust belt}
americano há um aspecto emblemático do fechamento de fábricas ligadas à
indústria automotiva, mas também de serralherias, fábricas de
eletrodomésticos e outros
\href{https://www.zotero.org/google-docs/?Nb0AyK}{(Cowie e Heathcott,
2003; High, 2003)}. Há uma lacuna nos estudos a respeito da composição
típica da indústria em cada território e seus impactos. Assim, não pode
ser excluída a possibilidade de fatores de caráter setorial
influenciarem as trajetórias observadas.

Também há uma notada ausência de discussão a respeito do setor de
construção e do setor ligado ao fornecimento de energia, água, gás e
saneamento -- em geral denominado \emph{utilities}. Um dos problemas
dessa ausência é justamente dificultar comparações setoriais e também
invisibilizar parte significativa da produção industrial.

Tabela 12: resumo dos casos de desindustrialização na literatura

\begin{longtable}[]{@{}
  >{\raggedright\arraybackslash}p{(\linewidth - 4\tabcolsep) * \real{0.2356}}
  >{\raggedright\arraybackslash}p{(\linewidth - 4\tabcolsep) * \real{0.1939}}
  >{\raggedright\arraybackslash}p{(\linewidth - 4\tabcolsep) * \real{0.5705}}@{}}
\toprule\noalign{}
\begin{minipage}[b]{\linewidth}\raggedright
\textbf{Local}
\end{minipage} & \begin{minipage}[b]{\linewidth}\raggedright
\textbf{Periodização}
\end{minipage} & \begin{minipage}[b]{\linewidth}\raggedright
\textbf{Características gerais}
\end{minipage} \\
\begin{minipage}[b]{\linewidth}\raggedright
América do Norte
\end{minipage} & \begin{minipage}[b]{\linewidth}\raggedright
1960-1990
\end{minipage} & \begin{minipage}[b]{\linewidth}\raggedright
Centrado na redução do emprego na manufatura, atenção à indústria
automotiva. Intensificação nos anos 1970 em diante com crise e mudanças
no cenário político nos anos 1980
\end{minipage} \\
\begin{minipage}[b]{\linewidth}\raggedright
Grã-Bretanha
\end{minipage} & \begin{minipage}[b]{\linewidth}\raggedright
1950-1990
\end{minipage} & \begin{minipage}[b]{\linewidth}\raggedright
Centrado na redução do emprego na indústria extrativa, especialmente
carvão e também na indústria pesada. Distinção dos impactos e aceleração
do processo na ascensão de discursos neoliberais do fim dos anos 1970 em
diante
\end{minipage} \\
\begin{minipage}[b]{\linewidth}\raggedright
Europa Continental
\end{minipage} & \begin{minipage}[b]{\linewidth}\raggedright
1950-2000
\end{minipage} & \begin{minipage}[b]{\linewidth}\raggedright
Redução do emprego no setor extrativo, especialmente carvão, e
siderurgia. Associado a mudanças tecnológicas e realocação de
investimentos. Maior gradualidade, especialmente no exemplo alemão
\end{minipage} \\
\begin{minipage}[b]{\linewidth}\raggedright
Leste Europeu
\end{minipage} & \begin{minipage}[b]{\linewidth}\raggedright
1990-2000
\end{minipage} & \begin{minipage}[b]{\linewidth}\raggedright
Redução do emprego industrial no setor de manufatura no contexto de
transição para economia de mercado nos anos 1990. Processo rápido e mais
intenso comparado a casos anteriores
\end{minipage} \\
\begin{minipage}[b]{\linewidth}\raggedright
Leste Asiático
\end{minipage} & \begin{minipage}[b]{\linewidth}\raggedright
1990-2010
\end{minipage} & \begin{minipage}[b]{\linewidth}\raggedright
Redução de emprego industrial no setor de manufatura no contexto de
mudanças organizacionais e no mercado de trabalho. Menos intenso e mais
mediado que no Leste Europeu
\end{minipage} \\
\begin{minipage}[b]{\linewidth}\raggedright
Brasil
\end{minipage} & \begin{minipage}[b]{\linewidth}\raggedright
1980-2020
\end{minipage} & \begin{minipage}[b]{\linewidth}\raggedright
Redução do emprego no setor de manufatura associado a crises econômicas
e abertura comercial. Crescimento econômico mais rápido do setor de
serviços nos anos 2000 em diante.
\end{minipage} \\
\midrule\noalign{}
\endhead
\bottomrule\noalign{}
\endlastfoot
\end{longtable}

Fonte: autoria própria

Pode-se, ainda, recortar a literatura a partir de seus enfoques quanto à
classe, a raça e o gênero. Uma clara conclusão é que os estudos são
majoritariamente baseados na perspectiva de trabalhadores e suas
comunidades, em parte pela formação do campo inicialmente calcada no
ativismo contra o fechamento de fábricas
\href{https://www.zotero.org/google-docs/?ZgRt6K}{(High, 2013, p.
995--996; Lawson, 2020, p. 2)}. Em particular, são enfatizados os
trabalhadores industriais afetados diretamente pelo fechamento de
fábricas ou minas, mas em alguns trabalhos grupos sociais urbanos com
interesses antagônicos as comunidades de trabalhadores, como classes
médias urbanas, podem aparecer, em geral nas discussões sobre projetos
relativos à memória ou reestruturação econômica de comunidades
\href{https://www.zotero.org/google-docs/?73nnBO}{(Hodson, 2019;
Hoekstra, 2020; Lewis, 2016)}.

Trabalhos voltados à dimensão racial parecem concentrados no exemplo
americano. Particularmente, preocupam-se com impactos maiores da
desindustrialização sobre as oportunidades futuras e qualidade de vida
da população negra
\href{https://www.zotero.org/google-docs/?vlrv6U}{(Berger e Engzell,
2022)}. Nas metrópoles americanas a desindustrialização tem também
ligação com uma gradual gentrificação, especialmente dos subúrbios, onde
frequentemente o fenômeno está associado à composição racial dos
territórios
\href{https://www.zotero.org/google-docs/?XohD0P}{(Highsmith, 2014;
Winant, 2019)}. Nas discussões centradas na Europa foi identificada
pouca discussão sobre como a questão racial modifica os impactos da
desindustrialização nos grupos sociais. Isso a despeito da importância
da população imigrante para reestruturação econômica das regiões como o
Ruhr \href{https://www.zotero.org/google-docs/?JhbIrM}{(Berger, Wicke e
Golombek, 2017)}.

Já na questão de gênero, alguns trabalhos abordam esta do ponto de vista
dos impactos da desindustrialização a respeito da masculinidade e dos
papéis de gênero
\href{https://www.zotero.org/google-docs/?iCkdQR}{(Nayak, 2006)}. Um dos
pontos centrais é como as mudanças nos processos de trabalho são vistas
como mudanças de formas de trabalho mais ``masculinas'' para trabalhos
mais ``femininos'' -- particularmente no setor de serviços. Estes se
refletem em requisitos profissionais distintos, frequentemente
estereotipados em termos de gênero, por exemplo a ideia de que
\emph{soft skills}\footnote{Para discussão sobre e caracterização das
  \emph{soft skills}, ver Thompson e Smith
  \href{https://www.zotero.org/google-docs/?BfUnet}{(2010, p. 99--103)}}
seriam mais tipicamente femininas. Há discussões sobre os impactos sobre
as mulheres de trabalhadores e, mesmo que em menor número, sobre
fechamento de fábricas onde a população empregada era majoritariamente
feminina \href{https://www.zotero.org/google-docs/?npTwIO}{(Clarke,
2011)}. Há pouca discussão sobre as relações entre desindustrialização e
o aumento da participação feminina na força de trabalho.

4. Causas, impactos e instituições mediadoras

A literatura examinada aponta como os processos de desindustrialização
estão ligados tanto às mudanças no modo de produção industrial que
produzem ganhos de produtividade, como às decisões de desinvestimento e
realocação produtiva das empresas
\href{https://www.zotero.org/google-docs/?F7H7gv}{(Kollmeyer, 2009,
2018)}. Entre os elementos chave, globalização
\href{https://www.zotero.org/google-docs/?6Of68o}{(Alderson, 1999;
Alderson e Nielsen, 2002)} e automação
\href{https://www.zotero.org/google-docs/?Wnb2hL}{(Berger e Engzell,
2022)} são mobilizados como causas parciais do fenômeno. No quesito
produtividade, em particular, há uma ênfase sobre como seu crescimento
acima da demanda reduziu a participação da indústria no PIB e estoque de
empregos \href{https://www.zotero.org/google-docs/?Q1XrCC}{(Rodrik,
2016)}.

Quando visto desse ângulo, a desindustrialização parece intimamente
ligada à viabilidade econômica da atividade industrial em um dado
território \href{https://www.zotero.org/google-docs/?YT2Gwk}{(Hoekstra,
2020, p. 697; Phillips, 2015, p. 558)}. Relaciona-se, assim, tanto à
possibilidade de realizar aquele trabalho com menor número de pessoas
como à facilidade em deslocar a atividade produtiva para territórios de
custos mais baixos, especialmente países mais pobres que observaram
ganhos de produtividade consideráveis
\href{https://www.zotero.org/google-docs/?W9IXhk}{(Tomlinson, 2016)}.
Isso implica numa relação com mudanças tecnológicas relativas à produção
e ao controle, inclusive financeiro. Essas mudanças têm base em avanços
em automação \href{https://www.zotero.org/google-docs/?gj4heC}{(Berger e
Engzell, 2022)} e em tecnologias de comunicação.

Longe de um determinismo tecnológico, entretanto, essas mudanças estão
ancoradas em decisões conscientes. O uso particular dado a essas
tecnologias está associado a um conjunto de estratégias gerenciais
voltadas à redução de custos, obtenção de benefícios fiscais e ambientes
regulatórios mais flexíveis, crescente financeirização e redução do
poder de barganha coletiva dos trabalhadores. Essas mudanças gerenciais
se configuram em estratégias de realocação produtiva, terceirização de
atividades e adoção de modelos de trabalho novos
\href{https://www.zotero.org/google-docs/?YGzUlj}{(Kollmeyer, 2009,
2018)}, como métodos \emph{lean}, produção \emph{just in time} e outras
formas de incrementar flexibilização do trabalho e produção.

Portanto, para além de uma discussão sobre redução de custos ou
produtividade, há uma discussão sobre a disputa pela captura de
valor\footnote{Para discussões sobre estratégia de captura de valor ver
  Henderson \emph{et al.}
  \href{https://www.zotero.org/google-docs/?FvfDr5}{(2011)}} produzido
advindo desses ganhos de produtividade. Bluestone e Harrison
\href{https://www.zotero.org/google-docs/?MdzJNx}{(1982)} associam esse
movimento a uma resposta corporativa à redução de suas taxas de lucro,
mas seria mais claro colocá-lo como uma resposta corporativa à captura
de valor por parte do trabalho\footnote{Para discussão sobre
  transformação nas estratégias de acumulação e o conflito entre capital
  e trabalhadores ver Harvey
  \href{https://www.zotero.org/google-docs/?eJL7Z3}{(1992)}}. Os efeitos
dessa disputa sobre o Estado seriam sentidos futuramente com a aplicação
de políticas neoliberais
\href{https://www.zotero.org/google-docs/?26mgUJ}{(Lawson, 2020;
Scheiring e King, 2023)}, mas a etapa inicial da disputa pela
apropriação do crescente valor produzido se deu na esfera institucional
dos mercados de trabalho.

É em função dessa dimensão política que o processo não pode ser visto
deslocado da trajetória histórica de cada território. Como enfatizado em
Bluestone e Harrison
\href{https://www.zotero.org/google-docs/?jPNmiv}{(1982)}, mas também
presente em Nettleingham
\href{https://www.zotero.org/google-docs/?fyB4VM}{(2019)}, Hodson
\href{https://www.zotero.org/google-docs/?G2DsO2}{(2019)} e Clarke
\href{https://www.zotero.org/google-docs/?u62AOc}{(2015)}, decisões
conscientes de desinvestimento e desregulação foram gradativamente
responsáveis pela perda de competitividade das indústrias nos
territórios desindustrializados. Portanto, para além de mera
racionalidade econômica, há uma necessidade de compreender essas
mudanças como uma transformação desejada advinda da disputa política.

Em resumo, para além da transformação tecnológica e da oportunidade de
negócios, houve intensa articulação política de forma a tornar a
desindustrialização viável e vantajosa. Isso permitiu uma crescente
captura do valor produzido nos territórios por uma elite econômica, em
detrimento da perda de poder econômico e político do trabalho
organizado, assim como de seus efeitos sobre as comunidades locais.

\hl{Essa análise afasta a possibilidade de uma leitura excessivamente
determinista do fenômeno, atribuindo-o como quase consequência natural
do crescimento da renda \emph{per capita}. Essa perspectiva foi típica
da economia neoclássica por muito tempo e inicialmente foi importada
para o repertório sociológico}
\href{https://www.zotero.org/google-docs/?gypx9o}{(Clark, 1957 e Bell,
1973, apud Kollmeyer, 2009, p. 1655--1645)}\hl{. A própria identificação
da ocorrência do fenômeno em graus de renda per capita muito distintos
ao longo do tempo}
\href{https://www.zotero.org/google-docs/?h6IPTX}{(Rodrik, 2016)}
\hl{aponta a existência de determinantes tecnológicos, organizacionais e
políticos para além da mera ``racionalidade econômica''. Prefere-se,
então, ler a correlação dos processos de desindustrialização com níveis
de renda \emph{per capita} como uma contingência histórica naquele
conjunto de economias nas quais essa mudança social se deu primeiro.}

Além de uma diversidade de causas, há uma diversidade de impactos. Como
já dito, o ritmo e temporalidade do processo parecem importar na
dimensão e característica de suas consequências sociais. Pode-se
destacar os impactos da desindustrialização no volume de empregos e, por
consequência, na renda e grau de proteção social desses empregos
\href{https://www.zotero.org/google-docs/?8tCeXm}{(Beatty e Fothergill,
2020)}, particularmente na medida em que postos de trabalho criados no
setor de serviço tendem a enfatizar uma dualização do mercado de
trabalho \href{https://www.zotero.org/google-docs/?pcGMyP}{(Lee, 2016,
p. 84--86)}.

Assim, é central o impacto da desindustrialização no crescimento da
desigualdade, especialmente em conjunto com o enfraquecimento de
sindicatos e outras organizações de trabalhadores
\href{https://www.zotero.org/google-docs/?cyxUuH}{(Kollmeyer, 2018)} e
na piora de diversos indicadores relativos ao bem-estar da população, a
exemplo da ocorrência de depressão, do suicídio, do uso de drogas e
mesmo do aumento de doenças cardíacas
\href{https://www.zotero.org/google-docs/?zmiS6A}{(Scheiring e King,
2023; Winant, 2019)}. Ademais, a perda de empregos tende a afetar
permanentemente a renda e a chance de ocupação de melhores postos de
trabalho pelas populações impactadas
\href{https://www.zotero.org/google-docs/?VjAxkf}{(Berger e Engzell,
2022)}, sendo os efeitos na renda mais fortes quando se observa a troca
de ocupação em indústrias tipicamente mais bem remuneradas para outras
de remuneração média inferior
\href{https://www.zotero.org/google-docs/?hD6Aqa}{(Bluestone e Harrison,
1982; Clark e Gibbs, 2020)}.

Há, ainda, já citadas consequências relativas à perda de identidades e
modos de vida \href{https://www.zotero.org/google-docs/?NKeJWc}{(Beatty
e Fothergill, 2020; Linkon, 2018; Mah, 2012)}. Outro resultado é a
gentrificação e o apagamento da memória social, relacionados aos
projetos de reorientação dos territórios desindustrializados
\href{https://www.zotero.org/google-docs/?3ItOId}{(Lewis, 2016)}. Muitos
trabalhadores sentem a perda de conexão com suas comunidades e senso de
pertencimento aos espaços públicos dos quais historicamente eles foram
parte \href{https://www.zotero.org/google-docs/?pQ8ete}{(Clarke, 2015;
Phillips, 2015)}.

Quanto ao Estado, são significativos os efeitos da desindustrialização
sobre a sua capacidade de arrecadação
\href{https://www.zotero.org/google-docs/?VRyXI9}{(Highsmith, 2014)} e
no incremento de gastos
\href{https://www.zotero.org/google-docs/?Dx7U8x}{(Beatty e Fothergill,
2020)}. Em particular, projetos de austeridade fiscal são comuns como
resposta à pressão econômica e social criada pela desindustrialização
\href{https://www.zotero.org/google-docs/?UQbXO1}{(Lawson, 2020)}. Para
além, há impactos sobre a percepção de populações sobre grupos políticos
que podem ser duradouras, em geral sendo enfatizada a capacidade do
grupo político no poder de responder ao fenômeno
\href{https://www.zotero.org/google-docs/?Hd7A5i}{(Phillips, Wright e
Tomlinson, 2019; Scheiring, 2020)}. Assim, estes processos parecem
interagir fortemente com transições políticas, seja como causa ou
consequência, e afetam por décadas as forças políticas organizadas em
torno dos territórios desindustrializados.

É possível, ainda, observar como a literatura trabalha a dimensão
relacional, que assume duas grandes ênfases: enquanto alguns trabalhos
destacam o processo como uma relação direcionada da atividade econômica
do setor secundário para o terciário
\href{https://www.zotero.org/google-docs/?g4rsPO}{(Beatty e Fothergill,
2020; Clark e Gibbs, 2020; Clarke, 2015; Kollmeyer, 2018; Phillips,
Wright e Tomlinson, 2019; Tomlinson, 2016, 2020; Winant, 2019)}; outros
a enfatizam como uma relação direcionada entre um território dito
``desindustrializado'' e novos territórios que passam por
industrialização
\href{https://www.zotero.org/google-docs/?QosUuZ}{(Clarke, 2015; Hodson,
2019; Kollmeyer, 2018; Strangleman, 2017; Tomlinson, 2020)}. Nesse
cenário é comum um influxo de imigrantes dos antigos territórios aos
novos \href{https://www.zotero.org/google-docs/?Efn8gu}{(Scheiring e
King, 2023; Winant, 2019)}, por vezes em modo pendular
\href{https://www.zotero.org/google-docs/?UOGMrF}{(Clark e Gibbs, 2020,
p. 43)}.

Há na literatura pouco debate sobre as decisões e estratégias dos
territórios para onde se deslocam as atividades industriais, em geral
tendo sua atuação estratégica reduzida à ideia de oferecer mercados de
trabalho com salários mais baixos ou incentivos de ordem fiscal e
regulatória
\href{https://www.zotero.org/google-docs/?7TG9zQ}{(Strangleman, 2017, p.
474)}. Isso destoa de achados da sociologia econômica sobre como a
construção social de novos territórios é mais complexa e planejada do
que aparenta, exigindo arranjos sociopolíticos particulares
\href{https://www.zotero.org/google-docs/?FEwbWu}{(Santos, 2007)}.

Essa baixa complexidade parece se refletir numa leitura relativamente
rasa da relação entre instituições e desindustrialização. Estados
nacionais são considerados importantes agentes, impactando ritmo e
efeitos sociais, especialmente via uso de mecanismos de seguridade
social \href{https://www.zotero.org/google-docs/?8dvTLj}{(Tomlinson,
2016)}, proteção do emprego
\href{https://www.zotero.org/google-docs/?OZsxYI}{(Phillips, 2015;
Phillips, Wright e Tomlinson, 2019)} e via liderança em projetos de
reestruturação econômica
\href{https://www.zotero.org/google-docs/?14S3Qj}{(Berger, Wicke e
Golombek, 2017)}. Entretanto, na maior parte dos trabalhos são
considerados monolíticos e não como conjunto de organizações ou
instituições com conflitos internos. Atores ligados ao Estado aparecem
apenas enquanto governos alternantes no nível do executivo federal ou
governos locais enquanto articuladores dos projetos de reestruturação
econômica, mas com pouca atenção ao papel das burocracias estatais,
empresas estatais, legislativo ou judiciário.

Sobre a sociedade civil destaca-se a participação e influência de
organizações ligadas à ação coletiva de trabalhadores, particularmente
sindicatos \href{https://www.zotero.org/google-docs/?nGvPUZ}{(Hodson,
2019; Kollmeyer, 2018; Phillips, 2015; Phillips, Wright e Tomlinson,
2019)}, vistos como capazes de, em parte, frear, estabelecer condições
melhores ou reduzir os impactos sobre as populações locais da perda de
empregos industriais. Outro conjunto de instituições relevantes são
organizações de moradores de certos territórios
\href{https://www.zotero.org/google-docs/?qfHjCP}{(Nettleingham, 2019)},
que são particularmente atuantes na gestão do território
desindustrializado e que buscam influenciar as estratégias voltadas a
reorientar o papel econômico dos territórios
\href{https://www.zotero.org/google-docs/?XYaD25}{(Hodson, 2019; Lewis,
2016)}.

Por fim, instituições de caráter econômico, em especial as empresas, são
colocadas como os elementos mais relevantes para o processo de
desindustrialização, sendo ele resultante primariamente de suas decisões
\href{https://www.zotero.org/google-docs/?9Cf4QP}{(Clarke, 2015;
Kollmeyer, 2018)}. Entretanto, apesar do enorme peso associado às suas
decisões, análises compreensivas sobre elas são quase ausentes na
literatura. Toma-se como premissa a ação racional e orientada a
interesses de acumulação por parte de empresas cujas decisões seriam
quase ``inevitáveis'' nessa perspectiva. Entretanto, algumas evidências
sugerem que essa relação não é tão simples. Alguns artigos apontam como
a relação entre empresas e territórios é mais complexa
\href{https://www.zotero.org/google-docs/?O7AVby}{(Highsmith, 2014)} e
que a decisão de desinvestimento está também ligada a disputas nas
economias políticas dos países
\href{https://www.zotero.org/google-docs/?bL2vpv}{(Bluestone e Harrison,
1982)}. Assim, parece precoce naturalizar as decisões de desinvestimento
e realocação produtiva como a única resposta racional ao contexto
apresentado. Assim, outra lacuna diz respeito à escassez de análises
acerca das perspectivas de empresas, elites econômicas ou trabalhadores
em atividade gerencial.

5. Discussão e conclusões

As discussões acima exemplificam como a desindustrialização não possui
uma definição consensual. Antes, possui características gerais
associadas às suas manifestações empíricas, mas estas não
necessariamente são amarradas em uma definição teórica coerente. Isso
torna possível autores divergirem até mesmo sobre a existência do
fenômeno em um dado caso.

Se não há consenso claro sobre o conceito, a maior parte da literatura
examinada compartilha de uma preocupação com a discussão dessa mudança
social do ponto de vista dos atores que o vivenciam mais diretamente,
particularmente daqueles entendidos como membros da classe trabalhadora
industrial, suas famílias e comunidades. São, portanto, tratados como
observadores privilegiados.

Entretanto, discussões de nível analítico macro são deixadas em segundo
plano, talvez com exceção dos trabalhos utilizando análise documental ou
estatística. Isso gera uma lacuna da compreensão sociológica a respeito
do envolvimento de Estado e empresas nas transformações. Assume-se muito
facilmente uma perspectiva baseada na escolha racional para os atores
econômicos, a qual Highsmith
\href{https://www.zotero.org/google-docs/?UHzgwr}{(2014)} mostra não ser
suficiente para compreender a formação das decisões de relocação
produtiva ou fechamento de fábricas.

Também acaba relegada ao segundo plano a discussão do poder e seu papel
na conformação do processo de desindustrialização, sendo as formas de
resistência dos trabalhadores quase sempre entendidas como insuficientes
para alterar significativamente o curso dessa transformação
\href{https://www.zotero.org/google-docs/?pGTrMe}{(High, 2022)}. Com
poucas exceções, como Berger, Wicke e Golombek
\href{https://www.zotero.org/google-docs/?9i2v0M}{(2017)}, há pouca
preocupação sobre a inter-relação entre movimentos de trabalhadores e
comunidades e sua capacidade efetiva de influenciar os destinos dos
territórios. Dada, também, a pouca frequência de trabalhos centrados no
Estado ou empresas, o poder acaba sendo uma categoria pouco
problematizada. As transformações impõem-se sobre as comunidades as
quais apenas reagem a elas, sem discussão dos interesses e estratégias
orientadores das decisões dos atores organizacionais envolvidos.

Apontar esse viés não implica defender uma perspectiva elitista ou um
foco excessivo em instituições formais. Antes disso, é justamente a
compreensão tanto do Estado quanto das elites econômicas enquanto grupos
de atores complexos, contraditórios e socialmente enraizados que permite
tornar estes acessíveis à ação coletiva da sociedade civil. Se Estado e
empresas não forem vistas em sua realidade, compreendidos em suas formas
de articulação e uso do poder, nunca será claro as maneiras como as
populações afetadas são capazes de reagir.

Esse quadro dialoga com análises como a de High
\href{https://www.zotero.org/google-docs/?RNRlkC}{(2022)}, que reflete
sobre a ação de trabalhadores ser muitas vezes ``infrutífera''.
Especialmente diante de uma perspectiva que deseja enfatizar as
contradições dessa transformação com esforços de neo ou
re-industrialização, como no caso brasileiro, é preciso tornar Estado e
empresas analiticamente ``vulneráveis'' aos contextos sociais. Isso
permite, por exemplo, ver como mudanças econômicas se refletem em novos
arranjos e aprendizados institucionais operando nos territórios
\href{https://www.zotero.org/google-docs/?chPHjZ}{(Ramalho, 2005)}.

Carece, portanto, de uma discussão mais aprofundada sobre como a ação
social desses atores coletivos dotados de uma dimensão de poder, Estado
e empresas em particular, é afetada pelas instituições nacionais e
regionais. A ausência dessa discussão é não compreender parte importante
do argumento de Bluestone e Harrison
\href{https://www.zotero.org/google-docs/?D50Tmf}{(1982, p. 15)} de que
``desindustrialização não apenas acontece''\footnote{Tradução própria}.
Isso implica discutir as estratégias de acumulação das organizações
econômicas e a ação do Estado na forma de suas políticas públicas e
organização burocrática.

Adicionalmente, se a bibliografia analisada permite uma boa leitura para
os casos das nações do Norte Atlântico, não é tão capaz de lidar por
completo com a realidade do fenômeno em territórios com trajetórias
históricas distintas, como o Brasil. Quando o processo de
desindustrialização ocorre simultaneamente com industrialização tardia,
\emph{neo} ou re-industrialização ou em comunidades nas quais nunca a
era industrial não teve o mesmo impacto, as perspectivas teóricas
oriundas de como esse processo é vivido nas nações onde ele primeiro se
evidenciou não apresentam o mesmo potencial explicativo. Dito de outra
maneira, a realidade é mais heterogênea do que uma leitura situada na
Europa Ocidental e América do Norte é capaz de dar conta.

É importante também salientar as lacunas e limitações desse próprio
texto. Primeiramente, é necessário ampliar de maneira sistemática o
escopo da análise para abarcar melhor trabalhos recentes das áreas de
economia, geografia e ciência política. Em segundo lugar, é preciso
utilizar critérios não sistemáticos para incluir na discussão mais
trabalhos a respeito dos casos fora do eixo Norte Atlântico -- com
ênfase no Brasil. Por fim, a falta de discussões mais aprofundadas sobre
o impacto da desindustrialização nas relações de gênero e raça, além de
maior reflexão sobre sua dimensão ambiental, são limitações deste
balanço. Essas lacunas apontam a necessidade de continuar, aprofundar e
complementar as discussões aqui presentes em trabalhos futuros.

\subsection{Desindustrialização no
Brasil}\label{desindustrializauxe7uxe3o-no-brasil}

Desafios da reindustrialização brasileira

Leonardo Nogueira Aucar - doutorando em Sociologia no PPGSA/UFRJ

\hl{Currículo reduzido: Doutorando em sociologia pelo programa de
pós-graduação em sociologia e antropologia (PPGSA) da UFRJ e pesquisador
do grupo Desenvolvimento, Ambiente e Trabalho (DTA). Estuda processos de
desindustrialização em perspectiva histórica-comparada. Suas principais
áreas de atuação são a sociologia histórica e a sociologia econômica.
Possui graduação em Comunicação Social pela PUC-Rio (2015) e mestrado em
sociologia pela UFRJ (2022).}

\textbf{Resumo:} esse texto discute os desafios às tentativas de
reindustrialização no Brasil atual, como a crescente intensidade de
capital da atividade industrial e a subsequente redução de sua
capacidade de geração de empregos. Discute o processo histórico de
desindustrialização pelo qual o país passa desde os anos 1980 e como
tais transformações afetam à articulação de políticas para o
desenvolvimento econômico centradas na indústria. Debate também novos
argumentos em favor de projetos de reindustrialização, elaborados no
contexto institucional da nova política industrial brasileira, como: a
questão do emprego ``de qualidade'', do incremento da capacidade de
crescimento econômico, de sua sustentabilidade ambiental e da autonomia
nacional em setores estratégicos. Por fim, é discutida a afinidade do
projeto de reindustrialização com novas demandas sociais, em especial
aquelas ligadas à redução da jornada de trabalho.

\textbf{Palavras-chave:} industrialização, desindustrialização, nova
indústria brasil, política industrial, desenvolvimento econômico

Introdução: breve contexto da desindustrialização brasileira

Os processos de desindustrialização se tornaram frequentes a partir da
segunda metade do século XX (Emery, 2019; High, 2013; Lawson, 2020;
Strangleman, Rhodes e Linkon, 2013). Estes podem ser mensurados pela
redução da participação do setor industrial, em especial da manufatura,
nos Produtos Internos Brutos (PIBs) dos países e em seu estoque de
empregos (Tregenna, 2009). Apesar dessa queda ter sido particularmente
intensa nas economias da Europa Ocidental e dos EUA, ela se manifestou
também em economias menos diversificadas e de menor PIB \emph{per
capita}, em especial na América Latina e Leste europeu (Rodrik, 2016).

Dentre todos esses casos, o brasileiro se destaca como um dos mais
intensos e duradouros. Com breves interrupções nessa trajetória, o
Brasil viu a participação da atividade industrial em seu produto bruto
se reduzir desde o início dos anos 1980 até o presente momento (Feijó e
Araujo, 2023; Morceiro, 2021; Oreiro e Feijó, 2010). Apenas o setor de
manufatura regrediu de 27\% em 1985 para 17\% nos dados de 1996 e para
11\% em 2019, quando considerada a série temporal corrigida por Morceiro
(2021).

Entre as causas do fenômeno a nível mundial a principal parece ser o
significativo incremento da produtividade do trabalho industrial, que
reduziu fortemente a demanda de trabalhadores (Rowthorn e Ramaswamy,
1999). Ao mesmo tempo, uma dinâmica de realocação produtiva,
especialmente em direção ao Leste Asiático e parte do Sul Global,
intensificou o fenômeno nas economias de industrialização mais antiga
(Kollmeyer, 2009).

Entretanto, tais fenômenos não podem ser compreendidos fora de seu
contexto histórico. Se nos países supracitados o início desse processo
ocorreu nos anos 1950, tal período para o Brasil ainda era de
industrialização e de predomínio de políticas desenvolvimentistas
ancoradas no papel do Estado (Perissinotto, 2023) e no capital
estrangeiro (Evans, 1979). Apenas em meados da década de 1970 a
estagnação do crescimento industrial brasileiro começa a exibir seus
sinais. Esta se dá no contexto da crise da dívida externa brasileira e
do crescimento da inflação (Hermann, 2015). Na medida em que o controle
inflacionário se torna a principal preocupação da política fiscal
nacional, a política industrial passa ao segundo plano (Cano, 2017),
fato evidenciado pelo abandono do terceiro Programa Nacional de
Desenvolvimento em fins da ditadura militar (Faria, 2016, p. 184--186).

A partir dos anos 1980 o país começa a passar por uma efetiva
desindustrialização, ao mesmo tempo que transaciona do regime militar
para a democracia. Em tal contexto de transformação das instituições
políticas, foi contínua a ênfase no combate à inflação como política
fiscal, especialmente após a eleição de Collor{[}1989-1992{]} em 1989,
quando se acelerou um processo de liberalização da economia e abertura
comercial (de Paula, 2022). Após sucessivos planos econômicos mal
sucedidos, o Plano Real consolidou em 1993 as bases de um novo modelo de
crescimento econômico baseado, entre outros pontos, no \emph{superavit}
fiscal, privatizações, juros altos e na paridade real-dólar, até a
flexibilização do câmbio em 1998 (Giambiagi, 2015; Pires, 2022). Modelo
esse que, se por um lado logrou o controle hiperinflacionário, não
conseguiu evitar a contínua desindustrialização. Em concomitância de tal
cenário com as reestruturações no modelo produtivo, a ascensão do
chamado Toyotismo (Ramalho, 2010), os anos 1990 marcaram significativa
queda no emprego industrial no Brasil (Ramalho e Rodrigues, 2013).

Aqui enfatiza-se, portanto, que não se trata de mero ``problema
econômico''. A desindustrialização brasileira precisa ser vista também,
para sua apropriada análise sociológica, como consequência do
direcionamento do capital político brasileiro desde os anos 1970 para a
estabilização inflacionária. Assim, enfatiza-se a capacidade de
articulação política para projetos de desenvolvimento como variável
importante para entendimento da trajetória da economia brasileira e para
análise de seus desdobramentos futuros.

Essa dinâmica acompanhou dois outros fenômenos agravantes.
Primeiramente, desde ao menos os anos 1990 o Brasil vive uma
desconcentração industrial, exibindo uma saída da atividade industrial
da região sudeste, em especial de São Paulo, para outras regiões
brasileiras (CNI, 2021). Tal fenômeno intensifica o impacto do processo
de desindustrialização em antigos territórios industriais, além de
frequentemente estar relacionado às disputas fiscais entre estados
(Arbix, 2000), via redução de impostos e oferta de subsídios, e ao
incremento da demanda por infraestrutura urbana (Bluestone e Harrison,
1982), reduzindo a capacidade de resposta fiscal (i.e. a oferta de
recursos financeiros disponíveis) do Estado ao fenômeno.

Em segundo lugar, o incremento do volume de empregos no chamado ``setor
de serviços'', que acompanhou a desindustrialização nos países de
industrialização avançada (Amsden, 2001), exibe características
particulares no Brasil. Isso devido, no caso nacional, a baixa
formalização, baixa remuneração, pouca qualificação e pouca perspectiva
de crescimento de carreiras nos principais ``novos'' setores
empregadores (Bridi, Oliveira e Salas, 2024). Com particularidades da
atividade econômica emergente em cada período, dos setores de
telemarketing dos anos 1990 até os entregadores de aplicativo nos dias
atuais , o incremento do uso de tecnologia e as inovações gerenciais não
se refletem num incremento na qualificação do trabalho, tornando o
Brasil espécie de paradigma \emph{``bravermaniano''} (Braverman, 1998)
da desqualificação do trabalho nos tempos contemporâneos.

Para além disso, o movimento acompanha uma tendência de precarização do
trabalho e dissolução das relações de trabalho formais (Ramalho, 2010).
No caso brasileiro esse fenômeno encontra evidência nas sucessivas
transformações institucionais (Thelen, 2003) que a legislação
trabalhista vem sofrendo (Ramalho, Salles Pereira Dos Santos e Jácome
Rodrigues, 2019). Tais processos de mudanças institucionais, assim como
as disputas fiscais mencionadas anteriormente, pressionam fiscalmente o
Estado brasileiro ao reduzir a sua capacidade de arrecadação.

Desafios da reindustrialização brasileira

Os dilemas fiscais, seja via aumentos dos gastos com seguridade social
ou seja via redução da arrecadação estatal pela redução do emprego
formal, impostos pelos processos de desindustrialização são um elemento
chave para compreender alguns dos principais desafios aos projetos de
reindustrialização. Como Winant (2019) aponta, essas pressões fiscais
tendem a ser refletidas em projetos de austeridade fiscal que, ao menos
parcialmente, buscam mitigar seu impacto. Assim, políticas de
austeridade não podem ser vistas apenas como resultado de uma concepção
ideológica arbitrariamente restritiva de projetos de desenvolvimento.
Elas precisam ser entendidas também como respostas a transformações
estruturais. Por sua vez, essas respostas são enviesadas pelos
interesses das frações de classe dominantes, obscurecendo respostas
alternativas de longo prazo. Isso não implica, entretanto, que
restrições fiscais a projetos de desenvolvimento, e em especial de
reindustrialização, não se imponham.

Dadas essas constrições, tanto as estruturais quanto as institucionais
relativas aos interesses da classe dominante, fica claro que a
capacidade de ação estatal possui limitações. Tais limitações são
relevantes na medida em que a ação do Estado é elemento necessário para
o desenvolvimento econômico, especialmente em economias de
desenvolvimento ou industrialização tardia e que se inserem numa
economia global diante de competidores já ``avançados'' (Amsden, 2001;
Gerschenkron, 1962). Especialmente quando se busca reorientar o perfil
das atividades econômicas de um país para além do que seriam suas
``vantagens competitivas'' presentes (Evans, 1995). Significativa
evidência dessa relação entre Estado e desenvolvimento já foi
apresentada para diversas nações (Evans, Rueschemeyer e Skocpol, 1985;
Rueschemeyer e Evans, 1985), inclusive o Brasil (Evans, 1979;
Perissinotto, 2023).

Para analisar tal dinâmica, primeiro pode-se analiticamente separar a
ação do estado na economia em dois eixos: ações voltadas a acrescer o
ritmo de acumulação e ações de caráter redistributivo (Rueschemeyer e
Evans, 1985). Dito de outra maneira, a ação do Estado precisa ser
pensada em sua capacidade de: 1- acelerar o ritmo de crescimento
econômico do país ou de setor econômico estratégico; 2- redistribuir os
ganhos econômicos entre diversos grupos sociais. Tal perspectiva é
vantajosa ao avançar de uma dicotomia de ``menos-mais'' Estado para
efetivamente pensar as formas específicas de intervenção estatal (Evans,
1995).

No caso brasileiro, tal dinâmica parece se refletir num \emph{trade-off}
entre esforços estatais para desenvolvimento, que podemos pensar como
estímulo ao investimento, e para ampliação da capacidade de consumo de
determinadas classes sociais. Não se trata aqui de uma oposição direta
entre consumo e investimento, mas de uma limitação na capacidade de ação
do Estado de promover ambos os objetivos \emph{sem incorrer em déficit}
ou tentando reduzi-lo ao máximo, dadas as consequências do déficit
público no modelo de desenvolvimento e ambiente institucional discutido
anteriormente -- pressão inflacionária e aumento de juros.

Essa análise parte da constatação de que, ao menos para o caso
brasileiro, a maior parte dos economistas de orientação ortodoxa como
também de algumas perspectivas heterodoxas influentes, como a novo
desenvolvimentista, consideram a existência de algum grau de equilíbrio
fiscal necessário ao processo de desenvolvimento (Bresser-Pereira,
2012). Assim, mesmo que essa política de superávit possa ser revertida
temporariamente, no que se denominaria uma política ``anticíclica'', a
dimensão fiscal tende a ser um fator limitante na capacidade de ação do
Estado no longo prazo.

A discussão acima serve para estabelecer o ponto de que o Estado não é
totalmente livre para livremente estimular investimento e consumo, sendo
necessário balancear em algum grau os dois objetivos. Como resultado se
encontra um problema político: qual consumo e em qual grau ele será
restringido para viabilizar o aumento da poupança e do investimento
produtivo? Como será direcionado o resultado do investimento da
poupança? Não há espaço aqui para reconstituir como, historicamente,
esse conflito se deu nos diversos modelos de desenvolvimento econômico
da história nacional -- por exemplo, o autoritarismo como forma de
ultrapassar a disputa política e do uso de dívida externa para contornar
a limitação fiscal durante a ditadura militar --, mas ele é tomado como
ponto de partida para expressar a complexidade do problema da
desindustrialização e das tentativas de reindustrialização.

Isso porque as mudanças associadas ao processo de desindustrialização
dificultam a resolução do problema político e econômico, que pudesse
propiciar um processo de reindustrialização. Notadamente, desde a
primeira revolução industrial e de forma acelerada a partir da metade do
século XX, a atividade industrial se tornou significativamente mais
complexa do ponto de vista tecnológico e mais intensa em capital
(Kollmeyer, 2009; Rowthorn e Ramaswamy, 1999), mesmo que com parcial
exceção de alguns casos no Leste Asiático. Tal processo é coerente com a
crescente dificuldade que economias têm de realizar o ``catch-up''
tecnológico ao longo do tempo (Amsden, 2001). Como já apontava
Gerschenkron (1962), cada ciclo de industrialização de um conjunto de
economias parece requerer maior capital e maior atuação do estado para
sua viabilização. Esse maior custo de capital financeiro também implica
maior necessidade de capital político a ser mobilizado, o que pode ser
desafiador quando as elites possuem interesses diretamente associados a
atividades econômicas dependentes do uso intensivo de trabalho, de
caráter extrativista ou rentista.

Ao mesmo tempo, essas características da atividade industrial propiciam
ao trabalho um grau maior de produtividade, especialmente após a segunda
metade do século XX, que implica menor demanda por trabalhadores
(Kollmeyer, 2009; Rowthorn e Ramaswamy, 1999). O resultado é que
processos recentes de industrialização, especialmente em setores com
alta perspectiva de crescimento, implicam na alocação de crescentes
volumes de capitais para decrescente número de empregos. A aliança entre
frações de classe que torna diversos projetos de industrialização
possíveis é fragilizada, especialmente em contextos democráticos. Na
ausência de uma elementos viáveis para mobilizar apoio popular e da
classe trabalhadora para avançar com projetos de reestruturação
econômica, o Estado tende a acomodar interesses mais conservadores de
suas elites, ancorados em vantagens competitivas já existentes.

Importante notar que, no passado nacional, essa ``acomodação'' envolveu
muitas vezes autoritarismo significativo, inclusive perante frações da
própria elite. Por exemplo, nos períodos Varguista e da Ditadura Militar
viabilizaram projetos de industrialização baseados na conciliação dos
interesses das elites (Perissinotto, 2023), inclusive as elites
internacionais no segundo caso (Evans, 1979). Mesmo assim, é importante
notar que em ambos os casos algum grau de apoio popular ou ao menos
apaziguação de demandas urbanas era preenchido pela noção do emprego
industrial. Na história brasileira talvez apenas o breve período dos
anos 1950 até o golpe de 1964 tenha conjugado de forma significativa
esse desenvolvimentismo com democracia (Barbosa, 2020). Um projeto de
reindustrialização por vias democráticas, se em parte possibilita maior
grau de promoção do bem-estar social e igualdade, ao mesmo tempo exige
maior presença e participação de trabalhadores, suas famílias e suas
instituições de classe para ser viabilizado.

A esse problema adiciona-se mais uma camada: a dificuldade de
coordenação de projetos de reindustrialização entre as diferentes
regiões do Brasil levando em conta suas particularidades. As já citadas
disputas fiscais entre estados são apenas um exemplo. Para além de
dificuldades de conciliar interesses, cabe ressaltar a dificuldade de
promover projetos de integração. Casos bem sucedidos de políticas
estatais para desenvolvimento econômico em geral contam também com a
presença de uma burocracia estatal extremamente capacitada, em algum
grau autônoma e bem integrada com a sociedade (Evans, 1995). Mesmo que
em alguns casos tal capacidade burocrática exista de forma limitada, ela
não aparenta ser em número suficiente e nem contar com o necessário
amparo político-financeiro e capilaridade para desenvolver e coordenar
projetos de nível local e regional na escala exigida (Evans, 1995).

Os novos argumentos pró-indústria

O cerne do argumento acima pode ser assim resumido: a reindustrialização
brasileira não é mero problema técnico-econômico que seja superável com
a correta combinação de taxa de juros, câmbio, subsídios públicos,
privatização, estatização ou qualquer instrumento econômico. Antes, ela
existe como um problema político de estabelecer um projeto nacional de
desenvolvimento ancorado na atividade industrial, que enfrenta
obstáculos por um lado devido às característica das elites nacionais e
por outro pela redução da capacidade da atividade industrial de agregar
apoio de fração significativa da classe trabalhadora devido a seu menor
potencial de geração de empregos.

Esse problema não passou ao largo das discussões sobre a
reindustrialização brasileira, mesmo que por vezes tenha sido feito
apenas de forma implícita. Como evidência, um conjunto de argumentos
``pró-industrialização'' tem se formado como resposta a esse fenômeno e
como tentativa de estruturar um projeto coerente de reindustrialização.
Alguns desses argumentos encontraram uma formalização institucional na
nova política industrial, adotada a partir de 2024, denominada Nova
Indústria Brasil (CNDI e MDIC, 2024). É a partir desse referencial que
discutirá-se em seguida alguns desses ``argumentos'', apontando
entretanto alguns de seus limites. Eles dizem respeito ao projeto de
reindustrialização como base de um crescimento econômico sustentável
gerador de empregos ``de qualidade'', como resposta à crise climática,
como ampliador da capacidade exportadora do país, como melhoria do
acesso da população urbana à infraestrutura e serviços de saúde e como
busca de soberania em setores estratégicos (CNDI e MDIC, 2024). Essa
análise fica clara quando observa-se o peso do ``Plano Mais Produção''
dentro do conjunto de instrumentos da nova política industrial, que
``concentra em seus quatro eixos - Mais Inovação, Mais Verde, Mais
Exportação, Mais Produtividade -- para o período de 2023 a 2026''
{[}...{]} ``R\$ 300 bilhões'' (CNDI e MDIC, 2024, p. 14). Desde seu
anúncio o valor total para o plano foi incrementado em outros 42,7
bilhões (Taiar, 2024).

Mas o que destaca-se aqui é menos o valor alocado ao plano, que ainda se
mostra muito inferior a outras iniciativas de apoio ao crescimento, como
os planos Safra (Agência Brasil, 2024), e mais as lógicas expressados no
que diz respeito aos retornos esperados à sociedade em função do
desenvolvimento industrial. Dá-se atenção a tais justificativas, pois
elas embasam a discussão sobre a necessidade de uma política pública
sobre o tema e, portanto, são articuladoras entre a ação do Estado e os
interesses dos diversos grupos sociais representados por este.

É preciso esclarecer que o objetivo deste texto não é uma análise
aprofundada da nova política industrial. Em que pese os limites dos
argumentos apresentados acima, o Nova Indústria Brasil representa um
importante esforço de institucionalização de uma política industrial
ancorada simultaneamente em processos de aprendizado institucional e no
impacto direto na vida da população. O que importa aqui, mais
especificamente, é como os objetivos e estratégias delineados no plano
representam não apenas um diagnóstico dos obstáculos ao desenvolvimento
industrial, mas associam a solução destes obstáculos diretamente com os
benefícios que possam ser compartilhados por diferentes grupos sociais.
É assim, um arranjo político que delineia certo conjunto de interesses
como prioritários e suficientes para justificar o uso de recursos
comuns.

Passando diretamente aos argumentos citados, o primeiro diz respeito à
necessidade de não apenas aumentar a produção e produtividade nacional,
mas de um impacto para a população na forma de ``empregos de qualidade''
(CNDI e MDIC, 2024). Em que pese a importância do incremento da
produtividade, sofisticação tecnológica e capacidade exportadora para o
crescimento do país, a articulação de um projeto em torno de ``empregos
de qualidade'' apresenta limitações. A primeira é a já citada limitação
na capacidade de produção de empregos em novas instalações industriais,
dado o paradigma da produção flexível (Ramalho, 2010). Uma segunda
problemática é que, mesmo em atividades mais ``sofisticadas'', a
qualidade do emprego não pode ser entendida como resultado necessário ou
mesmo esperado. Como mostra Ramalho (2006), novas unidades produtivas,
com práticas modernas e flexíveis, muitas vezes combinam essas práticas
com modelos de gestão do trabalho tradicionais e frequentemente
autoritários.

Isso implica que, ao passo que produtividade e sofisticação tecnológicas
são elementos necessários para existência de trabalhos de qualidade, não
são causas suficientes para eles (Mahoney, 2003). Exemplo recente, a
implementação da unidade produtiva de veículos elétricos em Camaçari tem
levantado embates diante das propostas salariais abaixo das expectativas
de trabalhadores e sindicatos (Gabriel, 2025). Tal discussão não é
surpreendente quando tomada em conta a literatura existente sobre a
relação entre processos de \emph{upgrade} econômico e \emph{upgrade}
social (Barrientos, Gereffi e Rossi, 2011). Na ausência de elementos
institucionais que garantam a captura do valor produzido por territórios
e comunidades -- como leis trabalhistas, um papel ativo do Estado ou
fortalecimento de organizações representativas dos trabalhadores --
processos de melhoria da produtividade podem não se refletir em
melhorias das condições sociais. No caso brasileiro parece inclusive
existir um movimento na direção contrária com processos de precarização,
através da flexibilização das leis trabalhistas, de redução da
seguridade social e enfraquecimento das organizações trabalhistas.
Assim, reduz a capacidade de mobilização política do emprego industrial
formal.

Tal discussão não significa apenas que projetos de reindustrialização
podem não vir a oferecer retornos sociais e econômicos esperados, mas
que a própria articulação do emprego industrial como base de apoio a
esses projetos fica fragilizada. Como justificar o significativo emprego
de capital econômico e político na manutenção de projetos de longo prazo
que apresentam retornos em empregos quantitativamente inferiores às
alternativas mais precárias e onde o retorno na qualidade desses
empregos encontra-se cada vez mais incerto?

Uma segunda é a questão ambiental e a proposta de posicionar a
reindustrialização como componente de uma agenda de transição
energética. Tal fator é enfatizado inclusive no lugar privilegiado que o
agronegócio tem na pauta do Nova Indústria Brasil (CNDI e MDIC, 2024).
Entretanto, assim como no caso do trabalho de qualidade, o arcabouço
institucional brasileiro parece direcionado de forma contra produtiva a
esse objetivo, vide recentes flexibilizações de regulações ambientais
(León, 2025). Nesse sentido, os esforços de construção de uma economia
sustentável e de longo prazo podem ser parcialmente anulados por uma
regulação ambiental que privilegia modelos econômicos menos sustentáveis
e de curto prazo.

Tanto no caso da busca por ganhos de produtividade industrial como na
incorporação do agronegócio via argumento de sustentabilidade, a
ampliação da capacidade de exportação é dado como objetivo privilegiado
(CNDI e MDIC, 2024). Isso é coerente do ponto de vista econômico, já que
a exportação foi importante ferramenta para desenvolvimento industrial
de países de industrialização tardia (Amsden, 2001), representando uma
forma de incremento da capacidade de acumulação, via entrada de
capitais, e também reposicionamento do país nas cadeias globais de valor
(Coe \emph{et al.}, 2004; Henderson \emph{et al.}, 2002). Entretanto, a
exportação implica, do ponto de vista da economia política, algumas
consequências.

Primeiramente, a orientação da economia às exportações implica, ao menos
no curto prazo, na contenção de consumo. Ao passo que o déficit em
balança comercial implica um grau de consumo adicional, um financiamento
de certo padrão de vida material, um superávit comercial tipicamente
implica um menor consumo doméstico. Isso fica notável ao observar
recentes impactos inflacionários da exportação de produtos alimentícios
(Moura, 2025) ou na perda de mercado de indústrias locais frente ao
acesso a produtos externos de baixo custo (Chmurzynski, 2024). Essa
problemática implica que, ainda que no longo prazo possa haver ganhos
com uma política de exportações, há no curto prazo uma disputa política
a respeito de qual consumo será inibido e qual classe, ou fração de
classe, arcará com tal custo.

Além disso, a busca de uma política de exportações tende a priorizar a
depreciação do câmbio, o que também é um elemento de restrição do
consumo. Nesse sentido, ela opera num sentido oposto de uma
sobrevalorização da taxa de câmbio, que permite manter padrões de
consumo elevados em troca de perda de competitividade nacional (Gala e
Libânio, 2011). Novamente, a problemática é menos técnica/econômica,
definir que um papel exportador é no longo prazo positivo a meta de
desenvolvimento econômico, e mais política, definir como a restrição do
consumo e seus supostos frutos futuros serão em cada caso repartidos
pelos grupos sociais.

Por fim, há de se ponderar o conjunto de argumentos voltados ao
desenvolvimento da infraestrutura urbana, além do estabelecimento de uma
indústria de saúde e defesa mais autossuficientes (CNDI e MDIC, 2024).
Em que pese a importância e a capacidade da melhoria da infraestrutura
urbana e da indústria de saúde em beneficiarem diretamente a qualidade
de vida da população, além do valor estratégico do setor de defesa e da
penetração do apelo ``nacionalista'' em parte da população, tais setores
apresentam limitações dentro do atual quadro fiscal brasileiro. Isso
pois investimentos em defesa e infraestrutura são altamente dependentes
da capacidade de gasto do Estado, de políticas públicas de longo prazo e
nem sempre resultam nos efeitos em cadeia desejados. Especialmente no
setor de defesa, expectativas de que a adoção de políticas de
desenvolvimento nacional resultem em ramificações e desenvolvimento de
atividades econômicas ``locais'' -- por exemplo via políticas de
Arranjos Produtivos Locais -- , tem tido resultados abaixo das
expectativas, mesmo que positivos, e dificuldades de continuidade em
cenários de mudança de governo (Ramalho e Conceição, 2024). São assim,
exemplo de como não basta o correto desenho de política pública, mas
também a necessidade de construção de um projeto de desenvolvimento
durável no tempo e resistente às alternâncias de poder democráticas.

Oportunidades da reindustrialização: um novo pacto baseado na
produtividade?

O cerne dos argumentos acima é a relação entre a existência de um
projeto político definido e a característica do processo de
desenvolvimento econômico. Isso implica o estabelecimento de uma relação
entre classes, ou frações de classes, que seja funcional a manutenção de
tal projeto político. As formas históricas que esse pacto assume variam
no tempo em acordo com os interesses e o poder -- a capacidade de
implementar seus interesses como interesses comuns (Dobbin, 2005) -- dos
diversos grupos sociais envolvidos.

No caso brasileiro o período de forte industrialização de 1967 até 1973
foi exemplo de um desses pactos. Como forma de sumarizar tal projeto,
poderia-se recorrer à frase do antigo ministro da economia Delfim Neto
de que é necessário ``\hl{fazer o bolo crescer, para depois dividi-lo}''
(Folha de São Paulo, {[}s.d.{]}). Essa visão expressada é a de um
desenvolvimento econômico acelerado, ancorado em dívida externa e papel
ativo do Estado, mas que requer significativo aumento da
desigualdade(Barbosa, 2022; Hermann, 2015). É uma ótima demonstração da
característica de um pacto firmado em ambiente de autoritarismo e onde
predominava a conciliação dos interesses de elites econômicas nacionais,
internacionais e da burocracia estatal (Evans, 1979). A consequência da
desarticulação desse projeto, em meio a crise da dívida externa e
inflação, foi a articulação de um projeto de desenvolvimento
alternativo, baseado em abertura comercial, privatização e, acima de
tudo, estabilidade inflacionária(Giambiagi, 2015; Paula, de, 2022).

Entretanto, se bem sucedido no combate inflacionário, tal projeto não
evitou a continuidade da desindustrialização nacional. Se os governos
Lula {[}2003-2010{]} e Dilma {[}2011-2015{]} em parte conseguiram uma
parcial repactuação desse modelo, com a inclusão da expansão da
capacidade de consumo das classes populares entre os interesses
privilegiados, especialmente via valorização salarial (Cosenza, 2022;
Fonseca, Arend e Guerrero, 2020), tão pouco conseguiu reverter uma
tendência à desindustrialização relativa -- crescimento menor da
indústria que outros setores -- que se prolongou nas últimas duas
décadas.

Em outras palavras, em meio ao uso de recursos políticos e econômicos em
projetos de redistribuição de renda bem sucedidos, o papel do Estado de
ampliar a capacidade de acumulação ficou em segundo plano. O resultado
foi um crescimento econômico robusto, baseado no incremento da demanda
pela classes média emergente, no bônus demográfico e no período de
\emph{boom} das \emph{commodities}, mas de pouca melhoria na
produtividade, pouco reposicionamento nas cadeias globais de valor e
recorrentes pressões inflacionárias (Cosenza, 2022; Fonseca, Arend e
Guerrero, 2020).

Também argumentou-se que, historicamente, os projetos de desenvolvimento
baseados em significativa industrialização dependem de crescentes
volumes de capital, ação estatal e se articularam em torno de dois
eixos: a aceleração da capacidade de crescimento econômico baseada em
incrementos da produtividade e a oferta significativa de empregos. Esses
eixos dizem respeito diretamente a interesses de frações da elite
econômica, tipicamente uma fração industrialista, e das camadas
populares.

Após hiato entre 2016 e 2022 na existência de uma política industrial
articulada e aposta num modelo liberal, pautado na redução dos custos do
trabalho, enxugamento do estado e exploração das ``vantagens
competitivas'', o país voltou a enfatizar a necessidade de um
desenvolvimento econômico que incorpore significativa reindustrialização
como Nova Indústria Brasil. Entretanto, a reestruturação da atividade
industrial em curso desde meados do século XX, que resultou em processos
de desindustrialização em diversos países, reduziu de forma
significativa a capacidade de geração de empregos dos investimentos
industriais. Em resposta a isso e também a pressões particulares do caso
brasileiro, a nova política industrial iniciada em 2024 parte da
articulação do argumento pró-indústria com outros interesses e demandas
da sociedade, a saber: uma melhoria na qualidade do emprego oferecido, a
possibilidade de um modelo de desenvolvimento mais ambientalmente
sustentável, a melhoria da infraestrutura urbana utilizada pela
população e a autonomia nacional em setores estratégicos, como saúde e
defesa.

Entretanto, mesmo sendo significativo avanço, tal projeto de
reindustrialização ainda encontra obstáculos. Alguns dos maiores têm
relação com a significativa constrição fiscal -- em parte decorrente da
baixa produtividade -- e da pouca capacidade de agregação de apoio
popular diante da baixa capacidade de gerar empregos e de um ambiente
institucional hostil a conversão de processos de \emph{upgrade}
econômico em \emph{upgrade} social. Nesse cenário, a continuidade de
esforços de incremento da capacidade de consumo da população,
especialmente via valorização real do salário mínimo, parece significar
uma priorização do papel redistributivo do Estado, mas encontra
resistência nas pressões inflacionárias, de causas em parte internas e
em parte externas. A resposta institucionalizada a essa pressão
inflacionária é um aumento significativo da taxa de juros, que reprime a
possibilidade de desenvolvimento industrial. Acende sinal de alerta a
estagnação do setor industrial nos dados do PIB do primeiro trimestre de
2025, com especial preocupação para a retração de 1\% da indústria de
transformação (Agência de Notícias - IBGE, 2025).

Daí a importância de se atentar cada vez mais ao conjunto de demandas
populares que pode constituir base de apoio político a novos projetos de
desenvolvimento. Esse elemento participativo é fundamental para
constituir um modelo de desenvolvimento que concilie as demandas
específicas da atividade industrial com as demandas de uma sociedade
democrática. Dentre muitas possibilidades, destaca-se aqui o
significativo impacto e popularidade de projetos relativos à redução da
jornada de trabalho, em particular a forte rejeição popular da chamada
``escala 6x1" (Martins, 2024). Se os incrementos de produtividade do
trabalho, inerentes e necessários ao desenvolvimento industrial, são
contraproducentes a projetos baseados na oferta quantitativa de
empregos, eles são aparentemente mais compatíveis com projetos baseados
no retorno dessa produtividade não apenas na forma de salários, e
consumo, mas na forma da redução do tempo de trabalho. Para isso,
entretanto, é preciso constituir um ambiente institucional que permita a
captação dessa produtividade pelos trabalhadores e não apenas pelo
capital, o que vai na contramão de um conjunto de reformas realizadas e
propostas nos últimos anos para (des)regulação do trabalho e capital.

A articulação dessa e de outras demandas sociais enquanto objetivos
plausíveis, para a população, de um projeto de reindustrialização é
precondição da mesma. Elas atuariam potencialmente como fonte de capital
político para um desenvolvimento econômico pautado em atividades mais
sofisticadas e intensivas em capital e menos ancorado naquelas
intensivas em trabalho. Parece, assim, constituir um dos principais
desafios e oportunidades na busca da reindustrialização brasileira.

\textbf{Bibliografia}

AGÊNCIA BRASIL. \textbf{Governo Federal lança Plano Safra 24/25 com R\$
400,59 bilhões para agricultura empresarial}. Disponível em:
\textless https://www.gov.br/planalto/pt-br/acompanhe-o-planalto/noticias/2024/07/governo-federal-lanca-plano-safra-24-25-com-r-400-59-bilhoes-para-agricultura-empresarial\textgreater.
Acesso em: 13 jun. 2025.

AGÊNCIA DE NOTÍCIAS - IBGE. \textbf{PIB cresce 1,4\% no primeiro
trimestre de 2025 \textbar{} Agência de Notícias}. Disponível em:
\textless https://agenciadenoticias.ibge.gov.br/agencia-sala-de-imprensa/2013-agencia-de-noticias/releases/43522-pib-cresce-1-4-no-primeiro-trimestre-de-2025\textgreater.
Acesso em: 13 jun. 2025.

AMSDEN, A. H. \textbf{The Rise of the Rest: Challenges to the West from
Late-Industrializing Economies}. Cary: Oxford University Press,
Incorporated, 2001.

ARBIX, G. Guerra fiscal e competição intermunicipal por novos
investimentos no setor automotivo brasileiro. \textbf{Dados}, v. 43, n.
1, p. 5--43, 2000.

BARBOSA, A. D. F. ``Developmentalist Brazil'' (1945-1964) as a concept:
historicizing and (re)periodizing development in Brazil.
\textbf{Brazilian Journal of Political Economy}, v. 40, n. 2, p.
332--354, jun. 2020.

BARBOSA, W. DO N. Alguns efeitos da política econômica durante a
Ditadura Militar (1964-1985). \emph{Em}: SOUZA, L. E. S. DE; PREVIDELLI,
M. DE F. S. DO C. (Eds.). . \textbf{História econômica do Brasil
contemporâneo}. Coleção Novos estudos de história econômica do Brasil.
Niterói, RJ\,: São Paulo, SP: Eduff\,; HUCITEC Editora, 2022. p.
243--298.

BARRIENTOS, S.; GEREFFI, G.; ROSSI, A. Economic and social upgrading in
global production networks: A new paradigm for a changing world.
\textbf{International Labour Review}, v. 150, n. 3--4, p. 319--340, dez.
2011.

BLUESTONE, B.; HARRISON, B. \textbf{The deindustrialization of America:
plant closings, community abandonment, and the dismantling of basic
industry}. New York: Basic Books, 1982.

BRAVERMAN, H. \textbf{Labor and monopoly capital: the degradation of
work in the twentieth century}. New York: Monthly Review Press, 1998.

BRESSER-PEREIRA, L. C. A taxa de câmbio no centro da teoria do
desenvolvimento. \textbf{Estudos Avançados}, v. 26, n. 75, p. 7--28,
2012.

BRIDI, M. A. D. C.; OLIVEIRA, R. V. D.; SALAS, M. M. Plataformas de
repartidores en América Latina: estado del arte. \textbf{Revista
Brasileira de Sociologia - RBS}, v. 11, n. 29, p. 14--40, 5 fev. 2024.

CANO, W. Brasil - construção e desconstrução do desenvolvimento.
\textbf{Economia e Sociedade}, v. 26, n. 2, p. 265--302, ago. 2017.

CHMURZYNSKI, G. \textbf{Substituição de produtos brasileiros por
importados prejudica indústria nacional}. Disponível em:
\textless https://noticias.portaldaindustria.com.br/noticias/economia/substituicao-de-produtos-brasileiros-por-importados-prejudica-industria-nacional/\textgreater.
Acesso em: 13 jun. 2025.

CNDI, C. N. DE D. I.; MDIC, M. DO D., Indústria, Comércio e Serviços.
\textbf{Nova indústria Brasil -- Nova indústria Brasil -- forte,
transformadora e sustentável\,: Plano de Ação para a neoindustrialização
2024-2026}. Brasília: CNDI, MDIC, 2024. Disponível em:
\textless https://www.gov.br/mdic/pt-br/composicao/se/cndi/plano-de-acao/nova-industria-brasil-plano-de-acao.pdf\textgreater.
Acesso em: 29 mar. 2024.

CNI, C. N. DA I. \textbf{Nota Economica 19 - Indústria fica menos
concentrada}: Nota Econômica. {[}s.l.{]} CNI, 2021.

COE, N. M. \emph{et al.} ``Globalizing'' regional development: a global
production networks perspective. \textbf{Transactions of the Institute
of British Geographers}, v. 29, n. 4, p. 468--484, dez. 2004.

COSENZA, A. C. Um brasil dos trabalhadores? Aspectos do desempenho
econômico brasileiro entre 2003 e 2016. \emph{Em}: SOUZA, L. E. S. DE;
PREVIDELLI, M. DE F. S. DO C. (Eds.). . \textbf{História econômica do
Brasil contemporâneo}. Coleção Novos estudos de história econômica do
Brasil. Niterói, RJ\,: São Paulo, SP: Eduff\,; HUCITEC Editora, 2022. .

DOBBIN, F. Comparative and Historical Approaches to Economic Sociology.
\emph{Em}: \textbf{The handbook of economic sociology}. 2nd ed ed.
Princeton, N.J.\,: New York: Princeton University Press\,; Russell Sage
Foundation, 2005. p. 26--48.

EMERY, J. Geographies of deindustrialization and the working‐class:
Industrial ruination, legacies, and affect. \textbf{Geography Compass},
v. 13, n. 2, p. e12417, fev. 2019.

EVANS, P. B. \textbf{Dependent Development: the Alliance of
Multinational, State, and Local Capital in Brazil}. Princeton, NJ:
Princeton University Press, 1979.

\_\_\_. \textbf{Embedded autonomy: states and industrial
transformation}. Princeton, N.J: Princeton University Press, 1995.

EVANS, P. B.; RUESCHEMEYER, D.; SKOCPOL, T. (EDS.). \textbf{Bringing the
state back in}. Cambridge {[}Cambridgeshire{]}\,; New York: Cambridge
University Press, 1985.

FARIA, C. E. D. \textbf{O papel da Confederação Nacional da Indústria na
política industrial brasileira (1938 -- 2014)}. {[}s.l.{]} Universidade
de Brasília, 5 dez. 2016.

FEIJÓ, C.; ARAUJO, E. (EDS.). \textbf{Industrialização e
desindustrialização no Brasil: Teorias, evidências e implicações de
política}. Curitiba, PR: Appris Editora e Livraria Eireli - ME, 2023.

FOLHA DE SÃO PAULO. \textbf{1968 - Ato Institucional 5 - Delfim Netto}.
Disponível em:
\textless https://www1.folha.uol.com.br/folha/treinamento/hotsites/ai5/personas/delfimNetto.html\textgreater.
Acesso em: 13 jun. 2025.

FONSECA, P. C. D.; AREND, M.; GUERRERO, G. A. Política econômica,
instituições e classes sociais: os governos do Partido dos Trabalhadores
no Brasil. \textbf{Economia e sociedade (São Paulo, Brazil)}, v. 29, n.
3, p. 779--809, 2020.

GABRIEL, D. \textbf{BYD oferece faixa salarial única de R\$ 1.950 e
irrita operáriosRepórter Brasil}, 21 maio 2025. Disponível em:
\textless https://reporterbrasil.org.br/2025/05/byd-faixa-salarial-unica-irrita-operarios-bahia/\textgreater.
Acesso em: 13 jun. 2025

GALA, P.; LIBÂNIO, G. Taxa de câmbio, poupança e produtividade: impactos
de curto e longo prazo. \textbf{Economia e Sociedade}, v. 20, n. 2, p.
229--242, ago 2011.

GERSCHENKRON, A. \textbf{Economic backwardness in historical
perspective: a book of essays}. Cambridge, Mass.: Belknap Pr, 1962.

GIAMBIAGI, F. Estabilização, Reformas e Desequilíbrios Macroeconômicos:
Os Anos FHC. \emph{Em}: VIANNA, S. B.; VILLELA, A. (Eds.). .
\textbf{Economia brasileira contemporânea}. {[}s.l.{]} Editora Atlas
Ltda, 2015. p. 166--195.

HENDERSON, J. \emph{et al.} Global production networks and the analysis
of economic development. \textbf{Review of International Political
Economy}, v. 9, n. 3, p. 436--464, 2002.

HERMANN, J. Reformas, endividamento externo e o ``milagre'' econômico
(1964-1973). \emph{Em}: VIANNA, S. B.; VILLELA, A. (Eds.). .
\textbf{Economia brasileira contemporânea}. {[}s.l.{]} Editora Atlas
Ltda, 2015. .

HIGH, S. ``The Wounds of Class'': A Historiographical Reflection on the
Study of Deindustrialization, 1973--2013. \textbf{History Compass}, v.
11, n. 11, p. 994--1007, nov. 2013.

KOLLMEYER, C. Explaining Deindustrialization: How Affluence,
Productivity Growth, and Globalization Diminish Manufacturing
Employment. \textbf{American Journal of Sociology}, v. 114, n. 6, p.
1644--1674, maio 2009.

LAWSON, C. Making sense of the ruins: The historiography of
deindustrialisation and its continued relevance in neoliberal times.
\textbf{History Compass}, v. 18, n. 8, p. e12619, ago. 2020.

LEÓN, L. P. \textbf{Comissões do Senado aprovam flexibilização do
licenciamento ambiental}. Disponível em:
\textless https://agenciabrasil.ebc.com.br/meio-ambiente/noticia/2025-05/comissoes-do-senado-aprovam-flexibilizacao-do-licenciamento-ambiental\textgreater.
Acesso em: 13 jun. 2025.

MAHONEY, J. Strategies of Causal Assessment in Comparative Historical
Analysis. \emph{Em}: MAHONEY, J.; RUESCHEMEYER, D. (Eds.). .
\textbf{Comparative historical analysis in the social sciences}.
Cambridge studies in comparative politics. Cambridge, U.K.\,; New York:
Cambridge University Press, 2003. p. 337--372.

MARTINS, L. \textbf{Datafolha: 64\% dos brasileiros querem o fim da
escala 6x1}. Disponível em:
\textless https://www.cnnbrasil.com.br/politica/datafolha-64-dos-brasileiros-querem-o-fim-da-escala-6x1/\textgreater.
Acesso em: 13 jun. 2025.

MORCEIRO, P. C. Influência metodológica na desindustrialização
brasileira. \textbf{Brazilian Journal of Political Economy}, v. 41, n.
4, p. 700--722, 2021.

MOURA, B. DE F. \textbf{Clima e foco em exportação explicam alta de
alimentos no longo prazo}. Disponível em:
\textless https://agenciabrasil.ebc.com.br/economia/noticia/2025-03/clima-e-foco-em-exportacao-explicam-alta-de-alimentos-no-longo-prazo\textgreater.
Acesso em: 13 jun. 2025.

OREIRO, J. L.; FEIJÓ, C. A. Desindustrialização: conceituação, causas,
efeitos e o caso brasileiro. \textbf{Revista de Economia Política}, v.
30, n. 2, p. 219--232, jun. 2010.

PAULA, R. Z. A. DE. Economia Brasileira na Nova República: 1985-1989.
\emph{Em}: SOUZA, L. E. S. DE; PREVIDELLI, M. DE F. S. DO C. (Eds.). .
\textbf{História econômica do Brasil contemporâneo}. Coleção Novos
estudos de história econômica do Brasil. Niterói, RJ\,: São Paulo, SP:
Eduff\,; HUCITEC Editora, 2022. p. 299--328.

PERISSINOTTO, R. \textbf{Ideias, burocracia e industrialização no Brasil
e na Argentina}. Rio de Janeiro, RJ: Eduerj, 2023.

PIRES, M. C. O modelo de mercado e a ``Segunda Década Perdida'':
1991-2000. \emph{Em}: SOUZA, L. E. S. DE; PREVIDELLI, M. DE F. S. DO C.
(Eds.). . \textbf{História econômica do Brasil contemporâneo}. Coleção
Novos estudos de história econômica do Brasil. Niterói, RJ\,: São Paulo,
SP: Eduff\,; HUCITEC Editora, 2022. p. 329--388.

RAMALHO, J. R. NOVAS FÁBRICAS, VELHAS PRÁTICAS: relações trabalhistas e
sindicais na indústria automobilística brasileira. \textbf{Caderno CRH},
v. 17, n. 41, 30 ago. 2006.

\_\_\_. Flexibilidade e crise do emprego industrial - sindicatos,
regiões e novas ações empresariais. \textbf{Sociologias}, v. 12, n. 25,
p. 252--284, dez. 2010.

RAMALHO, J. R.; CONCEIÇÃO, J. J. DA. \textbf{Caças supersônicos e o ABC
Paulista: Tecnologia e reconversão industrial}. São Paulo, SP: Editora
Papagaio, 2024.

RAMALHO, J. R.; RODRIGUES, I. J. SINDICATO, DESENVOLVIMENTO E TRABALHO:
crise econômica e ação política no ABC. \textbf{CADERNO CRH}, v. 26, n.
68, 2013.

RAMALHO, J. R.; SALLES PEREIRA DOS SANTOS, R.; JÁCOME RODRIGUES, I.
MUDANÇAS NA LEGISLAÇÃO TRABALHISTA, SINDICATO E EMPRESAS MULTINACIONAIS.
\textbf{Caderno CRH}, v. 32, n. 86, p. 343, 4 nov. 2019.

RODRIK, D. Premature deindustrialization. \textbf{Journal of Economic
Growth}, v. 21, n. 1, p. 1--33, mar. 2016.

ROWTHORN, R.; RAMASWAMY, R. Growth, Trade, and Deindustrialization.
\textbf{IMF Staff Papers}, v. 46, n. 1, p. 18--41, 1999.

RUESCHEMEYER, D.; EVANS, P. B. The State and Economic Transformation:
Toward an Analysis of the Conditions Underlying Effective Intervention.
\emph{Em}: EVANS, P. B.; RUESCHEMEYER, D.; SKOCPOL, T. (Eds.). .
\textbf{Bringing the state back in}. Cambridge {[}Cambridgeshire{]}\,;
New York: Cambridge University Press, 1985. p. 44--77.

STRANGLEMAN, T.; RHODES, J.; LINKON, S. Introduction to Crumbling
Cultures: Deindustrialization, Class, and Memory. \textbf{International
Labor and Working-Class History}, v. 84, p. 7--22, 2013.

TAIAR, E. \textbf{Governo amplia Nova Indústria Brasil para R\$ 342,7 bi
e define cadeias prioritárias e metas em saúde}. Disponível em:
\textless https://valor.globo.com/brasil/noticia/2024/08/15/governo-amplia-nova-indstria-brasil-para-r-3427-bi-e-define-cadeias-prioritrias-e-metas-em-sade.ghtml\textgreater.
Acesso em: 13 jun. 2025.

THELEN, K. How Institutions Evolve: Insights from Comparative Historical
Analysis. \emph{Em}: MAHONEY, J.; RUESCHEMEYER, D. (Eds.). .
\textbf{Comparative historical analysis in the social sciences}.
Cambridge studies in comparative politics. Cambridge, U.K.\,; New York:
Cambridge University Press, 2003. p. 208--240.

TREGENNA, F. Characterising deindustrialisation: An analysis of changes
in manufacturing employment and output internationally.
\textbf{Cambridge Journal of Economics}, v. 33, n. 3, p. 433--466, 1
maio 2009.

WINANT, G. ``Hard Times Make for Hard Arteries and Hard Livers'':
Deindustrialization, Biopolitics, and the Making of a New Working Class.
\textbf{Journal of Social History}, v. 53, n. 1, p. 107--132, 1 ago.
2019.
